\documentclass[11pt]{amsart}
\usepackage[T1]{fontenc}
\usepackage[utf8]{inputenc}
\usepackage{amsmath,amssymb,amsthm,mathtools}
\usepackage{microtype}
\usepackage[margin=1in]{geometry}
\usepackage[hidelinks]{hyperref}
\numberwithin{equation}{section}
\newtheorem{theorem}{Theorem}[section]
\newtheorem{proposition}[theorem]{Proposition}
\newtheorem{lemma}[theorem]{Lemma}
\newtheorem{corollary}[theorem]{Corollary}
\theoremstyle{definition}
\newtheorem{definition}[theorem]{Definition}
\theoremstyle{remark}
\newtheorem{remark}[theorem]{Remark}
\newcommand{\dd}{\,\mathrm d}
\newcommand{\E}{\mathbb E}
\newcommand{\Pp}{\mathbb P}
\newcommand{\R}{\mathbb R}
\newcommand{\C}{\mathbb C}
\newcommand{\one}{\mathbf 1}
\newcommand{\HH}{\mathsf H}
\newcommand{\OO}{\mathsf O}
\newcommand{\EE}{\mathsf E}
\newcommand{\BB}{\mathcal B}
\newcommand{\G}{\mathcal G}
\newcommand{\LL}{\mathcal L}
\newcommand{\resp}{\mathcal H}
\newcommand{\fb}{\mathfrak b}
\newcommand{\tf}{{\tau_{\mathrm f}}}
\newcommand{\cH}{c_{\mathrm H}}
\newcommand{\OmH}{\Omega_{\mathrm H}}
\newcommand{\norm}[1]{\left\|#1\right\|}
\newcommand{\abs}[1]{\left|#1\right|}
\newcommand{\ip}[2]{\left\langle #1,#2\right\rangle}
\newcommand{\out}{\mathrm{out}}
\newcommand{\per}{\mathrm{per}}

\hypersetup{
 pdftitle={Periodic probability densities for the full voltage-conductance equation},
 pdfauthor={Kaiwen Jin},
 pdfsubject={Periodic full-density solutions of a kinetic neuron equation},
 pdfkeywords={voltage-conductance, periodic solution, firing rate, reset boundary}
}

\title[Periodic voltage--conductance densities]{Periodic probability densities
for the full voltage--conductance equation}
\author{Kaiwen Jin}
\email{yui.ui.siki@gmail.com}
\date{}
\subjclass[2020]{35B10, 35Q84, 92B20}
\keywords{Voltage--conductance equation, periodic solution, firing rate,
reset boundary, kinetic Fokker--Planck equation}

\begin{document}

\begin{abstract}
We construct a nonnegative periodic probability density with nonconstant
firing rate for the full voltage--conductance neuron equation, for
suitable positive physical parameters. The model retains voltage
transport, conductance diffusion, threshold reset, reflection at zero
conductance, and firing feedback in both drift and diffusion.
A small-noise return-operator comparison and a corrected phase mode
give a genuine simple crossing of the full firing response.
Analytic parameter selection removes all other integer harmonic
resonances, and two free physical parameters close a fixed-period
rate equation. A positive prescribed period map then identifies
the analytic response with a nonnegative full density. The solution
is continuous in time into $L^1$, has the complete kinetic and
conormal boundary traces, finite entropy, and a uniformly finite
second conductance moment.
\end{abstract}

\maketitle
\section{Introduction}

The voltage--conductance equation describes a population of neurons
by a joint probability density in voltage and synaptic conductance.
Conductance diffusion supplies the noise, voltage transport carries
neurons to threshold, and a boundary reset returns a firing neuron
to its resting voltage. The firing flux also changes the conductance
drift and diffusion. The resulting autonomous equation can therefore
couple fluctuations of the population to its own firing rhythm.

Let $v_E>v_F>0$ and $y_L,y_*,a_*,c>0$. On
$D=(0,v_F)\times(0,\infty)$ put
\begin{align}\label{eq:full-coefficients}
 J(v,y)&=y(v_E-v)-y_Lv,&
 y_F&=\frac{y_Lv_F}{v_E-v_F},\notag\\
 K_F(t,y)&=y_*+cN(t)-y,& a_F(t)&=a_*+c^2N(t).
\end{align}
The firing rate is
$N(t)=\int_{y_F}^\infty J(v_F,y)F(t,v_F,y)\dd y$.
We seek a nonnegative mass-one density satisfying
\begin{equation}\label{eq:full-equation}
 \partial_tF+\partial_v(JF)+\partial_y(K_FF)-a_F\partial_{yy}F=0 .
\end{equation}
Writing $J_0(y)=J(0,y)$ and $J_F(y)=J(v_F,y)$, its voltage
traces obey
\begin{equation}\label{eq:full-voltage-boundaries}
 \begin{cases}
  F(t,0,y)=F(t,v_F,y)=0,&0<y<y_F,\\
  J_0(y)F(t,0,y)=J_F(y)F(t,v_F,y),&y>y_F.
 \end{cases}
\end{equation}
The conductance flux $Q=K_FF-a_FF_y$ has zero trace at $y=0$
and vanishes at infinity.
Voltage conditions are kinetic weak traces; the floor condition is
the natural conormal condition in the distributional formulation.
At infinity we mean that $Q$ paired against every smooth periodic
time test and smooth closed-voltage test has tail limit zero.
There is no prescribed value at the single grazing conductance $y_F$.

For this full equation, Fonte Sanchez and Mischler developed the
prescribed linear evolution and stationary theory, including
exponential relaxation in the frozen linear problem
\cite[Theorems 1.2 and 2.2]{FM24}. Fonte Sanchez, Mischler and Salort
proved global weak existence for general connectivity
\cite[Theorem 1.1 and Definition 4.1]{FMS26}. These results give
the evolution and trace framework used here. A nonconstant
periodic trajectory is a separate long-time question:
the voltage--conductance discussion in the survey of Carrillo
and Roux explicitly lists periodic-solution existence as open
\cite[Section 7.6]{CR25}. Carrillo, Dou and Zhou studied a
separable simplified equation and its fast-conductance limit;
their periodic solutions concern that limit
\cite[Section 4]{CDZ22}. The result below concerns the full
nonseparable transport field, positive conductance diffusion,
and both firing-feedback terms in \eqref{eq:full-coefficients}.

\begin{theorem}\label{thm:main}
There are positive constants $y_*,a_*,c,T$ such that, with
$v_E=2$, $v_F=1$ and $y_L=1$, equations
\eqref{eq:full-coefficients}--\eqref{eq:full-voltage-boundaries}
admit a nonnegative $T$-periodic distributional density $F$
of mass one whose firing rate $N$ is nonconstant.
The density is continuous in time into $L^1(D)$, has the
stated voltage and zero conductance-flux traces, and satisfies
\begin{equation}\label{eq:main-moment-entropy}
 \sup_{0\le t\le T}\int_D
       \bigl(y^2F(t,v,y)+|F(t,v,y)\log F(t,v,y)|\bigr)\dd v\dd y<\infty .
\end{equation}
One may choose
$F\in H^1_{\per}(0,T;L^2(D,e^{4y}\dd v\dd y))$ and
$N\in H^1_{\per}(0,T)$, with nonzero first Fourier coefficient
of $N$. The convention is $0\log0=0$.
\end{theorem}

The construction uses small frozen conductance diffusion and a mean
strictly above the deterministic firing threshold. One firing cycle
then provides a useful phase coordinate. We compare its noisy return
operator with an OU evolution over the deterministic cycle in a
common Gaussian weighted boundary space. The comparison holds in
operator norm for arbitrary data in that space. A corrected phase
mode, transferred through its isolated one-dimensional projection,
gives an actual simple characteristic root. Its real part has a
negative noise contribution and a positive mean-feedback contribution.
The reflecting-floor part of diffusion feedback is retained throughout.

Two further points are needed to obtain the theorem. First, the
fixed-noise physical resolvent is analytic in the mean and diffusion,
so a countable zero-set argument selects parameters with no other
integer harmonic resonance. The mean and coupling then serve as
two free parameters in a fixed-period rate equation. Second, its
initially signed density response is identified with the fixed
point of a positive prescribed period map. This step establishes
nonnegativity without requiring a pointwise positive lower bound
for the stationary density. Stability of the resulting periodic
solutions and explicit numerical parameter values are not determined
by this construction.

Section~\ref{sec:linear} constructs the divergence inverse with its
complete outgoing trace and natural floor input.
Section~\ref{sec:returns} establishes the process and cycle estimates.
Section~\ref{sec:crossing} proves the simple full-response crossing,
and Section~\ref{sec:periodic} closes the periodic rate equation and
proves the theorem. Appendix~\ref{app:boundary} gives the boundary
and process constructions; Appendix~\ref{app:clock} supplies the
bounded-clock comparison.

All complex transpose pairings below are bilinear; only energy
inner products use conjugation. We keep separate the random first
hit $\tf$, deterministic cycle time $\tau$, and final period $T$.
Constants in fixed-positive-noise resolvent estimates may depend on
that noise. Small-noise estimates are uniform only on the specified
compact mean and frequency sets, and harmonic expansions are for
each fixed harmonic.

\section{The linear equation and the firing trace}\label{sec:linear}

For the construction we fix
\begin{equation}\label{eq:physical-choice}
 v_E=2,\qquad v_F=1,\qquad y_L=1.
\end{equation}
Thus $D=(0,1)\times(0,\infty)$,
$J(v,y)=y(2-v)-v$, $y_F=1$, $J_0(y)=2y$ and $J_F(y)=y-1$.
We first keep the conductance mean and diffusion prescribed. Write them as
$\mu(t)$ and $\mathfrak a(t)$, bounded above and below by positive constants.
The conductance flux is
\[
 Q_{\mu,\mathfrak a}f=(\mu-y)f-\mathfrak a f_y.
\]
The weak formulation we use retains this flux at the reflecting boundary.
If $\gamma_0 f,\gamma_F f$ are the voltage traces, they satisfy
\begin{equation}\label{eq:voltage-traces}
 \gamma_0 f=\gamma_F f=0\quad(0<y<1),\qquad
 J_0\gamma_0 f=J_F\gamma_F f\quad(y>1).
\end{equation}
For a smooth test $\varphi$, compactly supported in $y$ and allowed to be
nonzero at all physical faces, the equation reads
\begin{align}\label{eq:weak-prescribed}
 \left[\int_D f\varphi\right]_{t_1}^{t_2}
 &+\int_{t_1}^{t_2}\!\int_1^\infty
 J_F\gamma_F f\,[\varphi(t,1,y)-\varphi(t,0,y)]\dd y\dd t\notag\\
 &=\int_{t_1}^{t_2}\!\int_D
 f\,[\varphi_t+J\varphi_v+(\mu-y)\varphi_y]
       -\mathfrak a f_y\varphi_y.
\end{align}
There is no boundary term at $y=0$. This is the zero conormal-flux condition.
For periodic solutions the endpoint term cancels when the test is periodic.
All voltage conditions are understood almost everywhere in their kinetic
trace measures. No trace value is needed at $y=1$.

\subsection{The frozen evolution}

Let $m$ belong to $I=[2,3]$ and let $a>0$. Fix $\sigma>1$ with
$2a\sigma<m$; we eventually take $\sigma=2$. Define
\[
 W(y)=e^{2\sigma y},\qquad
 \HH_\sigma=L^2(D,W\dd v\dd y),\qquad
 \HH_{-\sigma}=L^2(D,W^{-1}\dd v\dd y).
\]
\[
 \HH_{\sigma,0}=\left\{h\in\HH_\sigma:\int_Dh=0\right\}.
\]
We complexify these spaces. Their transpose pairing is the bilinear integral
$\langle h,u\rangle=\int_Dhu$; conjugation is used only in Hilbert energy
identities. The mass functional is bounded on $\HH_\sigma$.

We use the following consequences of the prescribed linear theory of
Fonte Sanchez and Mischler.

\begin{proposition}[The frozen linear evolution]\label{prop:frozen}
For prescribed positive bounded $\mu,\mathfrak a$ and initial data in
$\HH_\sigma$, there is a unique solution of the linear equation in
$C([0,t_*];\HH_\sigma)$, with local $L^2$ conductance derivative and the
renormalized weak traces, for every finite $t_*$. It preserves positivity
and mass. For $\mu=m$, $\mathfrak a=a$, its evolution is a strongly continuous
semigroup $S(t)$ on $\HH_\sigma$ with generator $\LL=\LL_{m,a}$.
There is a normalized nonnegative invariant density $M=M_{m,a}$, belonging
to every fixed exponential weighted space, and constants $C_a,\kappa_a>0$
such that
\begin{equation}\label{eq:frozen-gap}
 \norm{S(t)h}_{\HH_\sigma}\le C_a e^{-\kappa_a t}
           \norm{h}_{\HH_\sigma},\qquad h\in\HH_{\sigma,0}.
\end{equation}
The stationary density is in $D(\LL)$ and $\LL M=0$.
\end{proposition}

\begin{proof}
The prescribed statement and its renormalized Green identity are
\cite[Theorem 2.2 and (2.18)--(2.19)]{FM24}. The stationary and decay
statements follow from \cite[Theorem 1.2]{FM24} by setting the source's
connectivity to zero and its background mean and diffusion to $m,a$.
In particular its coefficients are then exactly $m-y,a$. The invariant
state is the generator-domain stationary state constructed in Section 7
of that paper. Autonomy and uniqueness give the semigroup property;
the asserted time continuity gives strong continuity. Its $p=2$ decay
estimate is \eqref{eq:frozen-gap}. These are estimates for the frozen
equation. The firing feedback will be included in Section~\ref{sec:crossing}.
\end{proof}

The domain $D(\LL)$ means the strong semigroup generator domain. We do not
replace it by a prescribed collection of smooth functions at the two
characteristic corners. The differential expressions in the interior are
\[
 \LL h=-\partial_v(Jh)-\partial_y((m-y)h)+ah_{yy},
 \qquad
 \G u=Ju_v+(m-y)u_y+au_{yy}.
\]

The outgoing space and the firing functional are
\begin{equation}\label{eq:out-space}
 \OO_\sigma=L^2((1,\infty),J_FW\dd y),\qquad
 \ell h=\int_1^\infty J_F(y)h_F(y)\dd y.
\end{equation}
The next estimate defines these quantities on the actual generator domain.

\begin{lemma}[Forward energy and the generator trace]\label{lem:forward-energy}
There is a constant $C=C_{m,a,\sigma}$ such that
\begin{align}\label{eq:forward-energy}
 \norm{S(t)h}_{\HH_\sigma}^2
 &+2a\int_0^t\norm{\partial_yS(s)h}_{\HH_\sigma}^2\dd s
 +\int_0^t\!\int_1^\infty J_F(1-J_F/J_0)
       |(S(s)h)_F|^2W\dd y\dd s\notag\\
 &\le \norm h_{\HH_\sigma}^2+
 C\int_0^t\norm{S(s)h}_{\HH_\sigma}^2\dd s.
\end{align}
For $h\in D(\LL)$, its derivative and outgoing trace satisfy
\begin{equation}\label{eq:domain-energy}
 2a\norm{h_y}_{\HH_\sigma}^2+\tfrac12\norm{h_F}_{\OO_\sigma}^2
 \le C\norm h_{\HH_\sigma}^2-2\operatorname{Re}(\LL h,h)_{\HH_\sigma}.
\end{equation}
In particular $\ell$ is continuous for the graph norm of $D(\LL)$.
The finite-time version of \eqref{eq:forward-energy} is uniform for
prescribed coefficients in a sufficiently small positive neighborhood of $m,a$.
\end{lemma}

\begin{proof}
Here are the truncations in the Green identity. For real $f$ use
\[
 b_n(s)=2n^2\bigl(\sqrt{1+s^2/n^2}-1\bigr),\qquad
 D_n(s)=s b_n'(s)-b_n(s).
\]
These are smooth convex functions with bounded first and second
derivatives, $0\le D_n\le b_n\le s^2$, and $b_n''\to2$. If boundedness of
the renormalization itself is required, first compose $b_n$ with a smooth
bounded argument truncation at height $A$. For fixed $n$ its first and
second derivatives are bounded independently of $A$. On compact physical
intervals the local $L^2$ derivative and the source trace bounds dominate
every term, so $A\to\infty$ recovers $b_n$. This also applies to the
entropy renormalizations used later.

Use $b_n$ with $W$ cut off between $R$ and $2R$. At high conductance the
reset condition, with $r=J_F/J_0\in(0,1)$, makes the boundary loss
\[
 J_F\bigl[b_n(f_F)-r^{-1}b_n(rf_F)\bigr]\ge0.
\]
It converges to $J_F(1-r)|f_F|^2$. Integrate the cross diffusion term as
$-a\int\partial_y b_n(f)W_y$. At the floor this gives
$-mD_n(f_0)+2a\sigma b_n(f_0)$. The source local $H^1_y$ class gives
$f_0\in L^2(\dd t\dd v)$, so this term converges to
$-(m-2a\sigma)|f_0|^2$. Its sign is used after, rather than before,
the truncation is removed.

Away from the cutoff the volume term is
\[
 (y+2)D_n(f)+2\sigma(m-y)b_n(f)+4a\sigma^2 b_n(f)
 \le C|f|^2,
\]
since $2\sigma>1$. The $y$-cutoff errors are bounded by a weighted
$\HH_\sigma$ tail: on their support $y/R$ is bounded. Pass $n\to\infty$
for fixed $R$, then $R\to\infty$. Fatou's lemma retains the nonnegative
gradient and reset losses. This proves \eqref{eq:forward-energy} and also
the favorable floor loss. Apply the argument to real and imaginary parts.

For $h\in D(\LL)$ divide by $t$. The norm difference quotient tends to
$2\operatorname{Re}(\LL h,h)$. The averages
$t^{-1}\int_0^t S(s)h\dd s$ tend to $h$ in graph norm. Jensen's inequality
and the preceding estimate bound their gradients and traces weakly;
their averaged Green identities identify the limits. As
$1-J_F/J_0\ge1/2$, weak lower semicontinuity gives
\eqref{eq:domain-energy}. Finally
\[
 |\ell h|\le
 \left(\int_1^\infty J_FW^{-1}\dd y\right)^{1/2}\norm{h_F}_{\OO_\sigma}.
\]
All coefficients and the limiting floor sign remain controlled in a
small positive prescribed-coefficient neighborhood.
\end{proof}

\subsection{Divergence inputs and the reflecting boundary}

The linearized feedback is a conductance divergence. For compact smooth
$g$ supported inside $D$, put $Bg=\partial_y g$. The completion for
general $g\in\HH_\sigma$ will be defined by
\begin{equation}\label{eq:B-definition}
 \langle Bg,\varphi\rangle=-\int_Dg\varphi_y,
\end{equation}
including tests nonzero at the floor. In particular it does not assign
a point value to $g(0)$.

By \eqref{eq:frozen-gap}, on the zero-mass space the ordinary inverse is
\begin{equation}\label{eq:ordinary-resolvent}
 R_{i\theta}h=\int_0^\infty e^{-i\theta t}S(t)h\dd t,
 \qquad \theta\in\R.
\end{equation}
It maps into $D(\LL)$ and $(i\theta-\LL)R_{i\theta}h=h$.
This follows by taking the difference quotient of the integral under
$S(s)$; the omitted integral over $(0,s)$ tends to $h$ after division by
$s$. The gap also excludes a homogeneous zero-mass solution.

We first bound divergence inputs independently of a killed backward
harmonic. The bounded-source backward construction needed in this proof
is given in Appendix~\ref{app:boundary}; in particular it proves the
actual Neumann condition before taking a global dual energy estimate.

\begin{proposition}[Completed divergence resolvent]\label{prop:divergence}
There is a bounded map
\[
 \mathcal R_\theta:g\longmapsto(h,h_y,h_F)
 \quad\hbox{from }\HH_\sigma\hbox{ to }
 \EE_\sigma:=\HH_{\sigma,0}\times\HH_\sigma\times\OO_\sigma,
\]
uniformly for all $\theta\in\R$, whose density component extends
$h=R_{i\theta}Bg$ for compact smooth $g$. Its weak equation is
\begin{align}\label{eq:completed-resolvent}
 i\theta\int_Dh\varphi
 -\int_Dh[J\varphi_v+(m-y)\varphi_y]
 +a\int_Dh_y\varphi_y
 &+\int_1^\infty J_Fh_F(\varphi_F-\varphi_0)\dd y\notag\\
 &=-\int_Dg\varphi_y.
\end{align}
Thus the natural flux $(m-y)h-ah_y-g$ has zero floor trace.
The original voltage conditions hold in the completed trace sense.
We continue to denote the density component by $R_{i\theta}Bg$.
The maps extend holomorphically to a neighborhood of each fixed
imaginary frequency where the ordinary zero-mass resolvent exists.
Constants and neighborhood sizes may depend on the fixed $a>0$.
\end{proposition}

\begin{proof}
Let $s$ exceed the full semigroup growth bound and a fixed dual-energy
constant. Here $Z_t=(V_t,Y_t)$ has reflected conductance
$\dd Y_t=(m-Y_t)\dd t+\sqrt{2a}\dd W_t+\dd K_t$ and voltage
$\dot V_t=J(V_t,Y_t)$; a threshold hit resets voltage from $1$ to $0$
without changing conductance. Proposition~\ref{prop:process} establishes
its density law, and Proposition~\ref{prop:backward-source} supplies
the transpose-resolvent representation below.
For bounded compact smooth $r$ the full-process value
\[
 u(z)=\E_z\int_0^\infty e^{-(s+i\theta)t}r(Z_t)\dd t
\]
is the transpose of the full forward resolvent. It is bounded by
$\norm r_\infty/s$. Proposition~\ref{prop:backward-source} in the appendix
gives, by localization at the corners and the actual Neumann floor,
\begin{equation}\label{eq:large-real-dual}
 (s-C)\norm u_{\HH_{-\sigma}}^2+
 a\norm{u_y}_{\HH_{-\sigma}}^2
 \le\norm r_{\HH_{-\sigma}}\norm u_{\HH_{-\sigma}}.
\end{equation}
Completion of compact smooth $r$, using closedness of the distributional
derivative, gives the same estimate for every $r\in\HH_{-\sigma}$.

For an arbitrary such $r$ set $r_0=r-(\int_DMr)\one$. The transpose of
\eqref{eq:ordinary-resolvent}, represented in $\int_D M\psi=0$, gives
\[
 (i\theta-\LL^T)\psi=r_0,\qquad
 \norm\psi_{\HH_{-\sigma}}\le C_a\norm r_{\HH_{-\sigma}}.
\]
The representative is an actual transpose-domain element: decompose
$h\in D(\LL)$ into $(\int h)M+h_0$. The stationary terms vanish by the
two mean-zero normalizations. The shifted identity is
\[
 \psi=(R^{\rm full}_{s+i\theta})^T(r_0+s\psi).
\]
Its input is bounded uniformly in $\theta$. Applying
\eqref{eq:large-real-dual} at this same frequency proves
$\norm{\psi_y}_{\HH_{-\sigma}}\le C_a\norm r_{\HH_{-\sigma}}$.

For compact $g$, $h=R_{i\theta}Bg\in D(\LL)$ and transpose duality now
gives
\[
 \langle h,r\rangle=\langle Bg,\psi\rangle=-\int_Dg\psi_y.
\]
Taking the dual norm bounds $\norm h_{\HH_\sigma}$ by $C_a\norm g_{\HH_\sigma}$.
Use \eqref{eq:domain-energy} with
$\LL h=i\theta h-Bg$. The remaining source term is
\[
 (Bg,h)_{\HH_\sigma}
   =-\int_Dg(\overline{h_y}+2\sigma\overline h)W.
\]
Young's inequality absorbs its derivative term. Hence
\begin{equation}\label{eq:triple-bound}
 \norm h_{\HH_\sigma}+\norm{h_y}_{\HH_\sigma}
       +\norm{h_F}_{\OO_\sigma}\le C_a\norm g_{\HH_\sigma},
\end{equation}
uniformly in $\theta$, including zero.

Approximate general $g$ by compact smooth functions away from $y=0$.
Apply \eqref{eq:triple-bound} to differences. All three components
converge strongly, and their weak equations give
\eqref{eq:completed-resolvent}. This is why the inhomogeneous floor flux
is retained even though the approximants vanish there. The product
$\EE_\sigma$ is an ambient Hilbert space; derivative and trace identities
are consequences of the weak equation for these outputs.

Finally, for $\lambda_0=i\theta$ use
\[
 R_\lambda B=R_{\lambda_0}B+
       (\lambda_0-\lambda)R_\lambda R_{\lambda_0}B.
\]
The second term has an ordinary $\HH_{\sigma,0}$ datum, so the ordinary
resolvent identity and domain energy give the local holomorphic extension.
\end{proof}

We will also need a finite-time form of the same completion.

\begin{lemma}[Causal divergence forcing]\label{lem:causal}
For $g\in L^2((0,t_*);\HH_\sigma)$ there is a unique zero-initial-data
solution of $(\partial_t-\LL)w=Bg$ in the completed energy class, with
\begin{equation}\label{eq:causal-bound}
 \sup_{0\le t\le t_*}\norm{w(t)}_{\HH_\sigma}^2+
 \int_0^{t_*}\bigl(\norm{w_y}_{\HH_\sigma}^2+
                  \norm{w_F}_{\OO_\sigma}^2\bigr)\dd t
 \le C_{a,m,t_*}\norm g_{L^2_t\HH_\sigma}^2.
\end{equation}
Its mass is zero and its natural base flux equals $g$ at the floor in
the weak formulation.
\end{lemma}

\begin{proof}
First take $g$ continuously differentiable in time and compactly supported
inside $D$. Its divergence is a differentiable $\HH_\sigma$ forcing, so
the semigroup convolution is a classical $D(\LL)$ solution. This also
follows by integrating the time derivative of $Bg$ in that convolution.
Apply \eqref{eq:domain-energy} to $w(t)$. Integrating the input as above,
Young's inequality and Gronwall give \eqref{eq:causal-bound}.
Complete the inputs in $L^2_t\HH_\sigma$. The outputs converge in
$C_t\HH_\sigma$, and their derivatives and traces converge in $L^2$.
Their weak equations retain the same floor input. The difference of two
completions is a homogeneous zero-initial-data solution, so prescribed
uniqueness in Proposition~\ref{prop:frozen} proves uniqueness. All these
solutions lie in its local energy class; a smaller exponential weight
may be used for this uniqueness comparison.
\end{proof}

\section{The reset process and one firing cycle}\label{sec:returns}

The first-return description will be used to locate a characteristic
root. We therefore identify both the state law and the firing-count
measure, rather than using a formal stochastic representation.

\subsection{The process, the first hit, and the outgoing law}

Let $Z_t=(V_t,Y_t)$, where
\begin{equation}\label{eq:reset-process}
 \dd Y_t=(m-Y_t)\dd t+\sqrt{2a}\dd W_t+\dd K_t,\qquad Y_t\ge0,
 \qquad \dot V_t=J(V_t,Y_t).
\end{equation}
Here $K$ is the minimal continuous regulator at zero. Whenever $V$ reaches
$1$ with $Y>1$, it is reset to zero without changing $Y$. Write $t_j$ for
these successive hits and $\tf$ for the next hit, including when the
initial voltage is not zero. The half-line Skorokhod map constructs $Y$
pathwise; the voltage equation then constructs the path between hits.
The absence of grazing and accumulation of hits is included in the next
proposition.

\begin{proposition}[Law and firing counts]\label{prop:process}
For a nonnegative $h\in\HH_\sigma$ of mass one, the process
\eqref{eq:reset-process} has density $S(t)h$ at every deterministic time.
Almost surely every threshold hit is outgoing and has $Y_{t_j}>1$.
There are finitely many hits on each bounded time interval, with finite
expected count. For every bounded Borel function $q$,
\begin{equation}\label{eq:count-measure}
 \E_h\sum_{t_j\le t}q(t_j,Y_{t_j})
 =\int_0^t\int_1^\infty q(s,y)J_F(y)(S(s)h)_F(y)\dd y\dd s.
\end{equation}
The same identity applies to integrable unbounded functions by
exhaustion.
\end{proposition}

The proof in Appendix~\ref{app:boundary} lifts the forward density to a
circle times a signed conductance line, applies Euclidean superposition,
and folds the conductance. A conditional-density argument excludes
grazing before identifying the circle path with the reset path.
Finally the voltage test $vq(t,y)$ identifies the entire count measure.
This proves the statement for the prescribed semigroup class in
Proposition~\ref{prop:frozen}; no converse for arbitrary weak solutions is
needed.

We now restrict $m$ to $I=[2,3]$ and decrease $a_0>0$ as necessary.
Set
\[
 \varepsilon=\sqrt a,\qquad \eta=\frac{y-m}{\varepsilon},\qquad
 \mathcal A=\partial_{\eta\eta}-\eta\partial_\eta,\qquad
 w_\beta(\eta)=e^{\beta\eta^2}.
\]
The reflecting floor is $\eta=-m/\varepsilon$.

\begin{lemma}[Uniform first-hit estimates]\label{lem:hitting}
There is $r_*>0$ such that
\begin{equation}\label{eq:hitting-time}
 \sup_{\substack{m\in I,\,0<a\le a_0\\z\in[0,1)\times[0,\infty)}}
                         \E_z e^{r_*\tf}<\infty.
\end{equation}
For $0<\beta<1/2$, after reducing $r_*$ if needed, every stopping time
$S\le\tf$ satisfies
\begin{align}
 \E_z e^{r_*S}w_\beta(\eta_S)&\le C_\beta w_\beta(\eta_0),
                                                   \label{eq:stopped-weight}\\
 \E_z\int_0^\tf e^{r_*s}w_\beta(\eta_s)\dd s
       &\le C_\beta w_\beta(\eta_0).                 \label{eq:killed-reward}
\end{align}
Constants are uniform on this small-noise family.
\end{lemma}

\begin{proof}
Choose $\delta>0$ with $4\delta<\min_I(m-1)$ and choose $t_0$ such that
$m(1-e^{-t_0})\ge1+2\delta$ on $I$. Couple $Y$ with the unreflected OU
process $U$ starting from zero. The explicit equations give
\[
 Y_t-U_t=e^{-t}Y_0+\int_0^t e^{-(t-s)}\dd K_s\ge0,\qquad
 U_t=m(1-e^{-t})+\varepsilon \zeta_t.
\]
For $B=t_0+1/\delta$, the stochastic-integral isometry and the maximal
inequality imply $\E\sup_{[0,B]}|\zeta_t|^2\le C_B$. For small $a$,
with probability at least $1/2$, $Y\ge1+\delta$ throughout $[t_0,B]$.
On this event, before a hit, $J(V,Y)\ge J(1,1+\delta)=\delta$.
A hit must occur by $B$, from every initial state. Independence of the
Brownian increments on successive blocks gives
$\Pp_z(\tf>nB)\le2^{-n}$ and hence \eqref{eq:hitting-time}.

For the weights,
\begin{equation}\label{eq:OU-weight}
 \mathcal A w_\beta
   =[2\beta-2\beta(1-2\beta)\eta^2]w_\beta.
\end{equation}
The regulator term is nonpositive because $w_\beta'$ is negative at the
floor. The expression $(\mathcal A+r_*)w_\beta$ is bounded above
uniformly in the floor location. Stop at a finite time, a finite upper
conductance and $S$. It\^o's formula and \eqref{eq:hitting-time} give
\eqref{eq:stopped-weight} by Fatou's lemma. The stronger inequality
$(\mathcal A+r_*)w_\beta\le C_\beta-w_\beta$, which follows from the
same negative quadratic term, gives \eqref{eq:killed-reward}.
\end{proof}

Let $N_0=\ell M$ be the stationary rate, and define the outgoing mark law
\begin{equation}\label{eq:mark-law}
 \pi_a(\dd y)=N_0^{-1}J_F(y)M_F(y)\dd y,\qquad y>1.
\end{equation}
We also write $\pi_a$ for its image under $y\mapsto(y-m)/\sqrt a$.
This law describes conductance at a firing time, whereas the conductance
marginal of $M$ describes a typical time.

\begin{lemma}[Stationary moments and cycle occupation]\label{lem:occupation}
The conductance marginal is exactly
\begin{equation}\label{eq:Gaussian-marginal}
 p_a(y)=C_a\exp[-(y-m)^2/(2a)]\,\one_{\{y\ge0\}}.
\end{equation}
For $L=\log2$, the stationary rate and mark law satisfy
\begin{equation}\label{eq:rate-mark-bounds}
 \frac{m-1}{L}\le N_0\le\frac{\E_M Y}{L},\qquad
 \int w_\beta\dd\pi_a\le C_\beta\quad(0<\beta<1/2).
\end{equation}
For every nonnegative measurable $g$, or every integrable complex $g$,
\begin{equation}\label{eq:occupation}
 \int_D Mg=N_0\,\E_{\pi_a}\int_0^\tf g(Z_s)\dd s.
\end{equation}
The cycle starts at voltage zero with conductance distributed as $\pi_a$.
In particular $N_0\E_{\pi_a}\tf=1$.
\end{lemma}

\begin{proof}
Integrating the stationary equation in voltage cancels the reset flux.
The marginal has zero current, so $ap_a'=(m-y)p_a$ and normalization
gives \eqref{eq:Gaussian-marginal}. In the $\eta$ variable it is a
standard Gaussian restricted to $\eta\ge-m/\varepsilon$, with
normalization at least $1/2$. Thus
\[
 \int_D M(1+|\eta|^j)w_\beta\le C_{\beta,j},\qquad \E_M Y\ge m.
\]
Put $X=\log[2/(2-v)]$ and $b(X)=e^X-1$. Between hits,
$\dot X=Y-b(X)$; a reset decreases $X$ by $L$. The stationary identity
with this bounded voltage test is
$LN_0=\E_M[Y-b(X)]$. Since $0\le b\le1$, it gives the first bounds
in \eqref{eq:rate-mark-bounds}.

To control the different measure $\pi_a$, use the concave truncations
$\zeta_n(s)=n(1-e^{-s/n})$ and test with
$(X/L)\zeta_n(w_\beta)$. The reset loss is $\zeta_n(w_\beta)$.
Its floor derivative is nonpositive, so the regulator has the favorable
sign. Concavity and \eqref{eq:OU-weight} imply
\[
 \G\bigl[(X/L)\zeta_n(w_\beta)\bigr]
       \le (Y/L+2\beta)w_\beta.
\]
Stationarity, followed by monotone convergence, bounds
$N_0\pi_a(w_\beta)$ by the preceding Gaussian moment. This proves the
second assertion in \eqref{eq:rate-mark-bounds}. Truncating in $y$ first
justifies this test in the forward identity.

For bounded $g$, put $k_g(z)=\E_z\int_0^\tf g(Z_s)\dd s$ and
$k_{g,0}(y)=k_g(0,y)$. The first-hit estimate makes $k_g$ bounded.
Dynamic programming, patched over the finitely many cycles, gives
\[
 \E_h[k_g(Z_t)-k_g(Z_0)]
   =-\E_h\int_0^t g(Z_s)\dd s+
                  \E_h\sum_{t_j\le t}k_{g,0}(Y_{t_j}).
\]
Start at $M$ and apply \eqref{eq:count-measure}. The left side is zero,
and division by $t$ proves \eqref{eq:occupation}. Monotone convergence
\cite[Theorem 1.4.44]{Tao11} and then domination give the asserted extensions.
\end{proof}

\subsection{Killed values and the generator pairing}

Fix henceforth $\beta=1/32$. The boundary space is
\[
 \BB_\beta=\left\{f\in C(\R;\C):
       f(\eta)/w_\beta(\eta)\longrightarrow0\text{ as }|\eta|\to\infty\right\},
 \qquad \norm f_\beta=\sup_\eta |f(\eta)|/w_\beta(\eta).
\]
Choose $r_0>0$ so that $4r_0<r_*$. For $|\operatorname{Re}\lambda|\le r_0$,
define
\begin{equation}\label{eq:killed-return}
 V_\lambda f(v,y)=\E_{v,y}[e^{-\lambda\tf}f(\eta_\tf)],
 \qquad
 P_a(\lambda)f(\eta)=V_\lambda f(0,m+\varepsilon\eta).
\end{equation}
For $\eta<-m/\varepsilon$ the last expression means its value at the
floor. Estimates \eqref{eq:stopped-weight}--\eqref{eq:killed-reward}
also apply when an interior reward is added. In particular a value with
terminal bound $F w_\beta$ and source bound $R w_\beta$ is bounded by
$C(F+R)w_\beta$. Discount derivatives follow from
$t^j e^{r_0t}\le j!(r_*-r_0)^{-j}e^{r_*t}$, giving local norm
holomorphy. Continuity of point-start values and the local backward
equation are proved in Appendix~\ref{app:boundary}.

The invariance of $\pi_a$ under $P_a(0)$ follows without an additional
renewal theorem. For a bounded terminal $f$, patch the harmonic
$V_0f$ over cycles and start at $M$. Its mean reset jump is
$\pi_a(f-P_a(0)f)$, hence is zero by \eqref{eq:count-measure}.
Monotone convergence extends this to nonnegative data, and domination
to $\pi_a$-integrable complex data. In particular it applies to $|f|^2$
for $f\in\BB_\beta$.

\begin{proposition}[A constant reset jump]\label{prop:constant-jump}
Suppose $u=V_\lambda f$ has the constant outgoing jump
\begin{equation}\label{eq:constant-jump}
 u(1,y)-u(0,y)=A,\qquad y>1.
\end{equation}
Then $u$ has at most polynomial growth at infinity, belongs to
$\HH_{-\sigma}$, and for every $h\in D(\LL)$,
\begin{equation}\label{eq:generator-pairing}
 \langle\LL h,u\rangle=\lambda\langle h,u\rangle-A\ell h.
\end{equation}
If $A=0$, this is the actual transpose-domain relation
$u\in D(\LL^T)$, $\LL^Tu=\lambda u$. At a frequency where the
zero-mass resolvent exists, the independently completed inverse in
Proposition~\ref{prop:divergence} gives
\begin{equation}\label{eq:divergence-pairing}
 u_y\in\HH_{-\sigma},\qquad
 A\ell R_\lambda Bg=-\int_Dg u_y,\quad g\in\HH_\sigma.
\end{equation}
\end{proposition}

\begin{proof}
Appendix~\ref{app:boundary} proves polynomial growth by stopping at a
fixed large conductance and estimating discounted counts before that
stop. It also constructs the local Neumann floor. Patch the killed
martingales over cycles to obtain
\[
 u(z)=\E_z[e^{-\lambda t}u(Z_t)]
                 +A\E_z\sum_{t_j\le t}e^{-\lambda t_j}.
\]
The required domination follows from polynomial growth, finite moments
of the reflected OU path supremum, and the pathwise bound
\begin{equation}\label{eq:count-bound}
 L\,\#\{j:t_j\le t\}\le X_0+\int_0^tY_s\dd s.
\end{equation}
Pair with $h$ and use \eqref{eq:count-measure}, splitting signed data
into positive densities when necessary:
\[
 \langle S(t)h,u\rangle=e^{\lambda t}\langle h,u\rangle
   -A\int_0^t e^{\lambda(t-s)}\ell S(s)h\dd s.
\]
For $h\in D(\LL)$, graph continuity of the trace in
Lemma~\ref{lem:forward-energy} permits differentiation at zero.
This proves \eqref{eq:generator-pairing}.

For compact smooth $g$ in the interior set $h=R_\lambda Bg$.
The resolvent equation and \eqref{eq:generator-pairing} imply
$A\ell h=\int_Du\,\partial_y g$. Proposition~\ref{prop:divergence}
bounds this functional by $C|A|\norm g_{\HH_\sigma}$.
Weighted $L^2$ Riesz representation therefore supplies the
distributional derivative $u_y\in\HH_{-\sigma}$. Complete the inputs
$g$ by the triple estimate \eqref{eq:triple-bound}; this gives
\eqref{eq:divergence-pairing}, including its natural floor flux.
The derivative of a killed harmonic was not used to construct that
inverse.
\end{proof}

\subsection{Cycle energies and the feedback functional}

For an interior source we use smooth functions on the closed physical
half-strip: on each finite part they admit smooth local extensions,
with the two voltage faces treated separately. No matching condition
on the source across reset is imposed. If $f\in\BB_\beta$ and
$|r(v,y)|\le Rw_\beta((y-m)/\sqrt a)$, define the inhomogeneous killed value by
\[
 u(z)=\E_z\left[e^{-\lambda\tf}f(\eta_\tf)
                    +\int_0^\tf e^{-\lambda s}r(Z_s)\dd s\right].
\]
The source is assumed smooth as just specified; a weighted size bound
alone is not a regularity assumption. The local construction in
Appendix~\ref{app:boundary} gives $\G u=\lambda u-r$ and the
reflecting derivative for these values. This is the source class
used in the next lemma and in its forced-feedback consequence.

\begin{lemma}[Stationary Fisher and killed-value energy]\label{lem:cycle-energy}
Uniformly for $m\in I$, $0<a\le a_0$,
\begin{equation}\label{eq:Fisher}
 a\int_D\frac{|M_y|^2}{M}\le \E_M Y+2\le C.
\end{equation}
The quotient is defined as zero where numerator and denominator vanish.
If $u$ is the killed value above with the stated smooth source,
terminal norm $F=\norm f_\beta$ and source bound $R w_\beta$, then
\begin{equation}\label{eq:cycle-energy}
 \int_D M|u|^2+a\int_D M|u_y|^2\le C(F+R)^2,
       \qquad |\operatorname{Re}\lambda|\le r_0.
\end{equation}
\end{lemma}

\begin{proof}
For the first estimate use the smooth convex entropy caps
\[
 b_{\delta,n}(0)=b_{\delta,n}'(0)=0,\qquad
 b_{\delta,n}''(s)=\frac1{(s+\delta)(1+s/n)^2},\quad s\ge0.
\]
Their even extensions are $C^2$; their first derivatives are bounded
for fixed $\delta,n$, and
$0\le s b_{\delta,n}'(s)-b_{\delta,n}(s)\le s$.
The bounded-argument extension in the proof of
Lemma~\ref{lem:forward-energy} makes these valid renormalizations.
Apply them to the stationary equation, first with a large-$y$ cutoff.
The floor loss and the reset loss are nonnegative by convexity.
For fixed $n,\delta$ the exponential tails of $M,M_y$ remove the
cutoff. Drop these losses, let $n\to\infty$ and then $\delta\downarrow0$.
The remaining volume coefficient is $y+2$, giving \eqref{eq:Fisher}
by Fatou's lemma. Interior smoothness and nonnegativity imply
$M_y=0$ at an interior zero of $M$; no boundary lower bound is used.

For the second estimate, the local equation and constructed floor permit
It\^o's formula for $|u|^2$ before $\tf$:
\[
 \G|u|^2=2\operatorname{Re}\lambda\,|u|^2
                 -2\operatorname{Re}(r\overline u)+2a|u_y|^2.
\]
Stop away from the grazing corners and at bounded conductance and time.
Point starts have no grazing hit, and a path cannot reach the floor
at voltage zero. The estimates with $4\beta$ and $4r_0$ give uniform
integrability for the endpoint terms; Fatou retains the gradient.
Start the cycle with $\pi_a$. Its invariance for $P_a(0)$, followed by
\eqref{eq:occupation}, yields
\[
 a\int_D M|u_y|^2\le
 \frac{N_0}{2}\pi_a|f|^2+
 |\operatorname{Re}\lambda|\int_D M|u|^2+
 \left|\int_DMr\overline u\right|.
\]
The reward bound, the exact Gaussian marginal, and
$\pi_a w_{2\beta}\le C$ bound every term on the right by $C(F+R)^2$.
They also bound $\int M|u|^2$, proving \eqref{eq:cycle-energy}.
\end{proof}

For a complex coupling $c$ define the complete weak feedback by
\begin{equation}\label{eq:feedback-functional}
 \fb_{\lambda,c}(f)
   =c\int_D M(V_\lambda f)_y-c^2\int_D M_y(V_\lambda f)_y.
\end{equation}
Its second term represents diffusion feedback, including the reflecting
floor. It is defined by first derivatives. For a smooth Neumann test
$u$ with integrable tails it equals $c^2\int M u_{yy}$.

\begin{corollary}[Small feedback and holomorphy]\label{cor:feedback}
On the preceding strip,
\begin{equation}\label{eq:feedback-bound}
 \norm{\fb_{\lambda,c}}_{\BB_\beta^*}
       \le C\left(\frac{|c|}{\sqrt a}+\frac{|c|^2}{a}\right).
\end{equation}
For a terminal-zero value with a smooth source in the same class and
source norm $R$, its feedback is bounded
by the same right side times $R$. The functional is norm holomorphic
in $\lambda$ and polynomial in $c$.
\end{corollary}

\begin{proof}
Cauchy--Schwarz with \eqref{eq:Fisher} and
\eqref{eq:cycle-energy} proves the bounds. For holomorphy, compact
interior tests satisfy
\[
 \int_D M\varphi(V_\lambda f)_y
       =-\int_DV_\lambda f\,\partial_y(M\varphi).
\]
The right side is holomorphic by the stopped estimates. Compact smooth
tests are dense in $L^2(M)$, and \eqref{eq:cycle-energy} bounds the
gradient locally uniformly. The pairing therefore extends
holomorphically to every $L^2(M)$ test. Take the tests $1$ and $M_y/M$,
the latter admissible by \eqref{eq:Fisher}. The uniform functional bound
then gives a norm-convergent Cauchy series. This argument requires no
positive lower bound for $M$ at a boundary.
\end{proof}

\section{A simple characteristic crossing}\label{sec:crossing}

We next turn a small-noise cycle calculation into a zero of the full
firing-response multiplier. The operator norm comparison and the
one-dimensional spectral projection are essential: a small residual
alone would not locate a characteristic root of this nonselfadjoint
problem.

\subsection{Phase coordinates and the return norm limit}

Recall $X=\log[2/(2-v)]$, $L=\log2$, and $b(X)=e^X-1$. Define
\begin{equation}\label{eq:phase}
 s(x)=\int_0^x\frac{\dd r}{m-b(r)},\quad
 \tau=s(L)=\frac1{m+1}\log\frac{2m}{m-1},\quad
 \omega=\frac{2\pi}{\tau},\quad Z(s)=\frac1{m-b(x(s))}.
\end{equation}
Here $\tau$ is the deterministic cycle duration; it differs from the
random first hit $\tf$ and from the nonlinear period $T$ used later.
Writing $A=m+1$, direct integration gives
\begin{equation}\label{eq:Z-formula}
 Z(s)=\frac1A+\frac1{mA}e^{As},\qquad
 \Delta Z=Z(\tau)-Z(0)>0.
\end{equation}
In $(s,\eta)$ variables the backward expression is exactly
\begin{equation}\label{eq:phase-generator}
 \G=(1+\varepsilon\eta Z(s))\partial_s+\mathcal A.
\end{equation}
Phase endpoints are identified by reset. Functions with equal endpoint
values have no reset jump; $Z$ itself has different one-sided values
there.

Let $Q_t$ be the OU operator with mean $e^{-t}\eta$ and variance
$1-e^{-2t}$.

\begin{proposition}[Return comparison in operator norm]\label{prop:return-norm}
For each compact set $\Lambda$ in a sufficiently small strip
$|\operatorname{Re}\lambda|<r_0$,
\begin{equation}\label{eq:return-norm}
 \sup_{m\in I,\,\lambda\in\Lambda}
 \norm{P_a(\lambda)-e^{-\lambda\tau(m)}Q_{\tau(m)}}_{\BB_\beta\to\BB_\beta}
       \longrightarrow0\qquad(a\downarrow0).
\end{equation}
The physical return operator, extended constantly below the reflecting
boundary as in \eqref{eq:killed-return}, maps $\BB_\beta$ into itself and is locally
norm holomorphic in $\lambda$. No rate in \eqref{eq:return-norm} is
asserted after the conductance cutoffs have been removed.
\end{proposition}

\begin{proof}
We give the three estimates and the order of their limits.
For $\kappa>0$,
\[
 \frac{\G(e^{\kappa X}w_\beta)}{e^{\kappa X}w_\beta}
 =\kappa(m+\varepsilon\eta-b(X))
                     +2\beta-2\beta(1-2\beta)\eta^2.
\]
After adding a fixed positive discount exponent this is nonpositive
outside a core $|\eta|\le R_\kappa$, uniformly in small $\varepsilon$.
Reflection again contributes nonpositively. Stop at the core entrance
$S$ or $\tf$. On $\tf<S$, $X_\tf=L$, so the terminal expectation gains
the factor $e^{-\kappa L}$. On $S<\tf$, the continued killed expectation
is bounded by $Cw_\beta(R_\kappa)$, and
$\E e^{rS}\one_{\{S<\tf\}}\le\E e^{r\tf}\le C$. Consequently
\begin{equation}\label{eq:return-rows}
 \sup_{|\eta|\ge M}
 \frac{|P_a(\lambda)f(\eta)|}{w_\beta(\eta)}
 \le\left[e^{-\kappa L}+Cw_\beta(R_\kappa)e^{-\beta M^2}\right]\norm f_\beta
\end{equation}
on physical rows. For the constantly extended rows, the estimate at the floor is
multiplied by $w_\beta(-m/\varepsilon)/w_\beta(\eta)\le1$.
First choose $\kappa$, then $M$. This makes the weighted row tails
small uniformly as $a\downarrow0$. The comparison operator has the same
property, since, with $v_t=1-e^{-2t}$,
\[
 Q_tw_\beta(\eta)=(1-2\beta v_t)^{-1/2}
     \exp\!\left(\frac{\beta e^{-2t}\eta^2}{1-2\beta v_t}\right).
\]
For $t=\tau(m)$ its exponent is strictly below $\beta$, uniformly on $I$.

Next fix $|\eta_0|\le M$ and $R>M$. For small enough $\varepsilon$,
the floor lies below $-R$ and the voltage speed before leaving
$|\eta|<R$ is bounded below by $j_*>0$. A hit or an exit therefore
occurs before $D=1/j_*$. The explicit unreflected OU representation
and the maximal inequality give
\[
 \Pp(\text{exit }(-R,R)\text{ before }\tf)\le \frac{C_D}{(R-M)^2}.
\]
By the stopped estimate with $2\beta$ and twice the discount exponent,
Cauchy--Schwarz bounds its terminal weighted contribution by a constant
depending on $M$ times the square root of this probability. The same
argument removes paths exiting $(-R,R)$ for $Q_\tau$. These bounds
hold for every datum in the unit ball of $\BB_\beta$.

On the fixed interval $(-R,R)$, phase is a positive clock. In reversed
phase time the killed evolution is
\begin{equation}\label{eq:clock-equation}
 q_\varepsilon(t,\eta)u_t=(\mathcal A-\lambda)u,\qquad
 q_\varepsilon(t,\eta)=1+\varepsilon\eta Z(\tau-t),\qquad
 u|_{\eta=\pm R}=0,\quad u(0)=f.
\end{equation}
Appendix~\ref{app:clock} constructs this evolution for arbitrary
continuous data and proves
\begin{equation}\label{eq:clock-sup}
 \norm{u_\varepsilon(\tau)-u_0(\tau)}_\infty
       \le C_{R,I,\Lambda}\varepsilon\norm f_\infty.
\end{equation}
Thus choose $\kappa,M$, then $R$, and finally $\varepsilon$. On the
fixed interval a unit $\BB_\beta$ datum has supremum at most
$e^{\beta R^2}$; \eqref{eq:clock-sup} now proves
\eqref{eq:return-norm}.

For fixed $a>0$, continuity of the return rows follows from continuous
Skorokhod and voltage dependence up to a transverse hit, point-start
non-grazing, and higher-weight uniform integrability. This includes
the floor, as proved in Appendix~\ref{app:boundary}. The constant extension
is continuous at that boundary. Estimate \eqref{eq:return-rows} gives vanishing weighted
tails, so the output belongs to $\BB_\beta$. The stopped discount
series gives norm holomorphy.
\end{proof}

\subsection{The corrected phase mode}

A stationary phase average is useful before constructing the mode.
For every continuous $\tau$-periodic function $h$, let
$\overline h=\tau^{-1}\int_0^\tau h(s)\dd s$ and take a periodic
primitive $q'=h-\overline h$. Along a stationary process $q$ has no
reset jump and no reflecting derivative. Therefore
\begin{equation}\label{eq:phase-average}
 \E_M h-\overline h
   =-\varepsilon\E_M[\eta Z(s)(h-\overline h)],\qquad
 |\E_M h-\overline h|\le C\varepsilon\norm h_\infty.
\end{equation}
The Gaussian marginal justifies the expectation. This identity avoids
any assertion of a joint phase--conductance product limit.

Fix a nonzero integer $k$ and set $q=k\omega$. Choose correctors with
equal endpoint values by
\begin{align}\label{eq:correctors}
 \phi'-\phi&=-iqZ,&
 \mathcal W&=Z(iq\phi+\phi'),&
 \lambda_2&=\overline{\mathcal W},\notag\\
 \psi'-2\psi&=-\mathcal W,&
 \chi'&=\lambda_2-\mathcal W,& \overline\chi&=0.
\end{align}
The first two periodic first-order equations are uniquely solvable,
and the last has a solution because its right side has mean zero.
All interval derivatives needed below are bounded for this fixed $k$,
uniformly in $m\in I$. Define initially
\begin{equation}\label{eq:base-mode}
 u_a=e^{iqs}\bigl[1+\varepsilon\eta\phi+
                  a((\eta^2-1)\psi+\chi)\bigr].
\end{equation}
Since $\mathcal A\eta=-\eta$ and
$\mathcal A(\eta^2-1)=-2(\eta^2-1)$, the coefficients of orders
$1,\varepsilon,a$ in $(\G-iq-a\lambda_2)u_a$ vanish exactly.
Using a smooth cutoff, make the corrections flat in a fixed-width layer at the physical
floor $\eta=-m/\varepsilon$. They remain unchanged above that layer.
The resulting $u_a$ has exact Neumann derivative at the floor and zero
reset jump. Cutoff errors are polynomial in $\varepsilon^{-1}$ times
a Gaussian-small weighted tail, hence are exponentially small.

Let the feedback of this smooth field be
\[
 \alpha=c\int_D M(u_a)_y-c^2\int_D M_y(u_a)_y.
\]
Put
\begin{equation}\label{eq:approx-mode}
 u_{\rm app}=u_a+\alpha(s/\tau)e^{iqs},\qquad
 \lambda_{\rm app}=iq+a\lambda_2+\alpha/\tau,\qquad
 f_{\rm app}=(u_{\rm app})_F.
\end{equation}
The added term has zero $y$ derivative. The reset jump and the feedback
of $u_{\rm app}$ both equal $\alpha$.

\begin{lemma}[Damping, gain, and residual]\label{lem:quasimode}
For a fixed $k\ne0$ and a fixed $C_0$, uniformly for complex
$|c|\le C_0a$,
\begin{align}
 \norm{f_{\rm app}-\one}_\beta&\le C\sqrt a,&
 \sup_{\text{physical }s,\eta}
       \frac{|(\G-\lambda_{\rm app})u_{\rm app}|}{w_\beta(\eta)}
       &\le Ca^{3/2}.                             \label{eq:mode-residual}
\end{align}
There are real $D=D(m)>0$ and complex $G_k=G_k(m)$ such that
\begin{align}
 \operatorname{Re}\lambda_2&=-k^2D,&
 \frac{\alpha}{\tau}&=cG_k+O(|c|\sqrt a+|c|^2),     \label{eq:mode-gain}\\
 G_k&=\frac{\Delta Z}{\tau^2}
           \frac{iq}{(A+iq)(1+iq)},&
 \gamma_k:=\operatorname{Re}G_k
   &=\frac{\Delta Z}{\tau^2}
       \frac{q^2(A+1)}{(A^2+q^2)(1+q^2)}>0.       \label{eq:gamma}
\end{align}
Moreover $\gamma_k/k^2$ strictly decreases for positive integers $k$.
\end{lemma}

\begin{proof}
Let the real periodic $\mathcal C$ solve $\mathcal C'-\mathcal C=-Z$.
Then $\phi=iq\mathcal C$, and
\[
 \operatorname{Re}\lambda_2
   =-q^2\overline{Z\mathcal C}
   =-q^2\overline{\mathcal C^2}=-k^2D,\qquad
 D=\omega^2\overline{\mathcal C^2}>0.
\]
Here $Z=\mathcal C-\mathcal C'$ and periodicity give the second equality.
Positivity follows also from
$\mathcal C(s)=\int_0^\infty e^{-t}Z(s+t\ {\rm mod}\ \tau)\dd t>0$.

The leading conductance derivative of \eqref{eq:base-mode} is
$e^{iqs}\phi(s)$, continuous at the phase seam. Formula
\eqref{eq:phase-average} replaces its $M$ average by its phase average
with $O(\sqrt a)$ error. The next derivative term is $\sqrt a$ times
$\eta$ times a bounded phase function; the floor error is Gaussian
small. With the Fourier convention
$Z_j=\tau^{-1}\int_0^\tau Z(s)e^{-ij\omega s}\dd s$,
\[
 Z_{-k}=\frac{\Delta Z}{\tau(A+iq)},\qquad
 \phi_{-k}=\frac{iqZ_{-k}}{1+iq}.
\]
Integration by parts in the diffusion term is valid for the smooth
Neumann field and bounds $\int M(u_a)_{yy}$ uniformly. These facts
give \eqref{eq:mode-gain}. Taking the real part yields
\eqref{eq:gamma}; its displayed denominators prove the strict decrease.

In particular $\alpha=O(a)$. The exact cancellations in
\eqref{eq:correctors} leave $O(a^{3/2})$ times a fixed polynomial in
$\eta$. The added term in \eqref{eq:approx-mode} contributes
$O(\alpha\varepsilon)+O(\alpha a)$ after its leading frequency shift
has canceled. All polynomial terms are bounded in the weighted
supremum norm. The floor cutoff contributes exponentially small errors.
Extending $f_{\rm app}$ constantly below the reflecting boundary preserves
the reset identity.
These observations prove \eqref{eq:mode-residual}.
\end{proof}

\subsection{An actual root and the full response}

Define the coupled return operator
\begin{equation}\label{eq:coupled-return}
 K(\lambda,c)=P_a(\lambda)+\one\otimes\fb_{\lambda,c}
                 \quad\hbox{on }\BB_\beta.
\end{equation}
By Corollary~\ref{cor:feedback} and Proposition~\ref{prop:return-norm},
for $|c|\le C_0a$ it converges in norm to
$e^{-\lambda\tau}Q_\tau$ on fixed complex neighborhoods of
$ik\omega$.

\begin{proposition}[The simple return root]\label{prop:return-root}
For each fixed $k\ne0$ and all sufficiently small $a$, there is a unique
root $\lambda_k$ near $ik\omega$ of the simple eigenvalue-one equation
for $K$. It is simple and holomorphic in $c$ for $|c|<C_0a$, and
\begin{align}
 \lambda_k&=ik\omega+a\lambda_2+cG_k+O(a^{3/2}),   \label{eq:true-root}\\
 \partial_c\lambda_k&=G_k+O(\sqrt a)              \label{eq:root-derivative}
\end{align}
on interior $c$ disks of radius comparable to $a$. Constants are
uniform in $m\in I$ for fixed $k$. There are real positive
$\cH(a,m)$ and $\OmH(a,m)$ satisfying
\begin{equation}\label{eq:crossing}
 \lambda_1(\cH)=i\OmH,\qquad
 \frac{\cH}{a}=\frac D{\gamma_1}+O(\sqrt a),\qquad
 \OmH\to\omega,\qquad
 \operatorname{Re}\partial_c\lambda_1(\cH)>0.
\end{equation}
\end{proposition}

\begin{proof}
Let $\Pi_\gamma f=(\int f\dd\gamma)\one$, where $\gamma$ is standard
Gaussian measure. The Mehler formula gives
\begin{equation}\label{eq:Mehler-gap}
 \norm{Q_{n\tau}-\Pi_\gamma}_{\BB_\beta\to\BB_\beta}\le Ce^{-n\tau}.
\end{equation}
For completeness, differentiate its Gaussian kernel in the mean
$e^{-t}\eta$ and variance $1-e^{-2t}$. After multiplication by
$e^{-\beta\eta^2}$, the kernel difference is bounded, for large $t$,
by $Ce^{-t}$ times a polynomial in $\eta$ times
$e^{-\beta\eta^2/2}$. Integrating against $w_\beta$ proves the
bound; the finitely many earlier times are absorbed in $C$.
The estimate is uniform for $\tau(m)$ on its compact interval.

It follows that the constant eigenvalue of $Q_\tau$ is simple and that
the complement resolvent outside $|\zeta|=e^{-\tau}$ is given by a
uniformly convergent series. Choose a small circle around $1$ in the
$\zeta$ plane and a sufficiently small fixed $\lambda$ disk around
$ik\omega$. The limiting resolvent on that circle is uniformly
bounded. The actual resolvent exists there by a Neumann series.
Its contour projection $\Pi=\Pi_{a,c,\lambda}$ tends in norm to
$\Pi_\gamma$ and has rank one: projections at distance less than one
are injective on each other's ranges. Resolvent identities give
idempotence and commutation. Normalizing by Gaussian mean gives a
holomorphic eigenvalue $e(a,c,\lambda)$ with
\[
 e\longrightarrow e^{-\lambda\tau},\qquad
 e_\lambda\longrightarrow-\tau
\]
at the harmonic; derivative convergence on a smaller disk follows
from Cauchy's formula.

At $\lambda_{\rm app}$ let
$v=V_{\lambda_{\rm app}}f_{\rm app}$ and $w=u_{\rm app}-v$.
It has zero outgoing terminal value and source
$r=-(\G-\lambda_{\rm app})u_{\rm app}$ bounded by
$Ca^{3/2}w_\beta$. This source is smooth on the closed physical
half-strip: the phase correctors are smooth up to their separate
voltage endpoints, and the floor cutoff is smooth.
Thus it belongs to the source class of Lemma~\ref{lem:cycle-energy}.
The killed-reward bound and the forced version of
\eqref{eq:feedback-bound} therefore give
\[
 (I-K(\lambda_{\rm app},c))f_{\rm app}
       =w(0,\cdot)+\fb_{\rm field}(w)\one,\qquad
 \norm{(I-K)f_{\rm app}}_\beta\le Ca^{3/2}.
\]
Here $\fb_{\rm field}(w)=c\int Mw_y-c^2\int M_yw_y$.
Since $\Pi f_{\rm app}\to\one$, projecting this identity yields
$|e(\lambda_{\rm app},c)-1|\le Ca^{3/2}$.

On a smaller fixed convex disk choose
$|e_\lambda+\tau|<\tau/3$; $e_{\lambda\lambda}$ is bounded there.
With $b=e_\lambda(\lambda_{\rm app},c)$, the map
\[
 \delta\longmapsto
 \delta-\frac{e(\lambda_{\rm app}+\delta,c)-1}{b}
\]
is a contraction on $|\delta|\le C'a^{3/2}$ for sufficiently large
fixed $C'$. Its value at zero is $O(a^{3/2})$, and its Lipschitz
constant on that disk is $O(a^{3/2})$. It gives an actual root with
\eqref{eq:true-root}. If two roots lay in the fixed convex disk,
integrating $e_\lambda+\tau$ on the segment between them would give
$\tau|\lambda-\lambda'|\le(\tau/3)|\lambda-\lambda'|$.
Thus the root is unique there and simple.

Holomorphy and Cauchy's formula on $c$ disks of radius comparable to
$a$ turn the uniform $O(a^{3/2})$ error into the $O(\sqrt a)$ derivative
error in \eqref{eq:root-derivative}. Choose
$C_0>4\sup_I D/\gamma_1$. On a real interior interval of $c$ of
length $C_1a$, the real part of $\lambda_1$ starts negative, ends
positive, and has strictly positive derivative. Its unique zero
gives \eqref{eq:crossing}. This concerns the indicated branch;
no stability assertion about the remaining spectrum is required.
\end{proof}

We now compare this root with the physical forward response. Define
\begin{equation}\label{eq:physical-response}
 g_c=-cM+c^2M_y,\qquad
 \resp(\lambda;m,a,c)=\ell R_\lambda Bg_c .
\end{equation}
The response is defined on the entire imaginary axis and locally
holomorphic near each frequency by Proposition~\ref{prop:divergence}.

\begin{proposition}[Response crossing and the static block]\label{prop:response-crossing}
At the crossing in \eqref{eq:crossing},
\begin{equation}\label{eq:response-root}
 \resp(i\OmH;m,a,\cH)=1,\qquad
 \partial_\lambda(1-\resp)(i\OmH;m,a,\cH)\ne0.
\end{equation}
For $c=O(a)$,
\begin{equation}\label{eq:static-response}
 \resp(0;m,a,c)=O(\sqrt a).
\end{equation}
\end{proposition}

\begin{proof}
Near the chosen nonzero harmonic $I-P_a(\lambda)$ is Fredholm of index
zero. Indeed the norm limit splits into a one-dimensional part and a
complement on which $I-P$ is invertible; block elimination leaves a
finite-dimensional block. On the nonzero imaginary axis its kernel
is zero. If $(I-P)f=0$, then $u=V_\lambda f$ has zero reset jump.
Proposition~\ref{prop:constant-jump} puts $u$ in the actual transpose
domain. Its $M$ mean is zero since $\lambda\ne0$, and the frozen gap
excludes this imaginary eigenfunction. Thus $I-P$ is invertible at
the crossing and in a possibly $a$-dependent complex neighborhood.

Set $f_0=(I-P)^{-1}\one$. Its killed value has jump one. The completed
Green pairing, with both feedback terms, gives
\begin{equation}\label{eq:response-match}
 \resp(\lambda)=\fb_{\lambda,c}(f_0),\qquad
 (I-K)f_0=(1-\resp(\lambda))\one.
\end{equation}
At a $K$ root on this imaginary axis the jump of an eigenfunction
cannot be zero, by the same frozen transpose argument. The first
identity follows from \eqref{eq:divergence-pairing}; the second
therefore implies $\resp=1$.

Let $\nu$ be the left eigenfunctional of the simple eigenvalue of $K$,
normalized by $\nu(f_0)=1$ at the root. It satisfies
$\nu(I-P)=\nu(\one)\fb$. If $\nu(\one)=0$, invertibility of $I-P$
would force $\nu=0$. Differentiating \eqref{eq:response-match} and
pairing with $\nu$ thus gives
\[
 -e_\lambda=\partial_\lambda(1-\resp)\,\nu(\one).
\]
This proves \eqref{eq:response-root}, including simplicity.

At zero frequency, $P_a(0)\one=\one$. Its simple spectral projection
is $\Pi_a f=\pi_a(f)\one$, by the invariant mark law. The inverse of
$I-P_a(0)$ on $\pi_a$-mean-zero data is uniformly bounded, by the
same separated complement. Let $t(z)=\E_z\tf$, with reset value $t_0$.
The occupation formula gives $\pi_a(1-N_0t_0)=0$. Solve
\[
 (I-P_a(0))f=1-N_0t_0,\qquad \pi_a f=0,
 \qquad w=V_0f-N_0t.
\]
Here $t$ has zero terminal datum and source $r=1$, while $w$ has
source $r=-N_0$ at $\lambda=0$. These constant sources satisfy the
smoothness and weighted bounds in Lemma~\ref{lem:cycle-energy}.
Then $\G w=N_0$ and $w_F-w_0=1$. Killed-reward and forced cycle
energy estimates give $\int M|w|^2+a\int M|w_y|^2\le C$.
The polynomial-growth argument in Appendix~\ref{app:boundary}
also applies with this bounded interior source. The generator
pairing now has the constant volume term $N_0\int h$. It vanishes
for $h=R_0Bg_c$, whose mass is zero. The same Riesz completion
therefore gives
$\resp(0)=c\int Mw_y-c^2\int M_yw_y=O(\sqrt a)$.
For small $a$, the static multiplier $1-\resp(0)$ is nonzero.
\end{proof}

\section{Periodic rates and nonnegative full densities}\label{sec:periodic}

The selected first characteristic mode does not by itself produce a
periodic density. We first remove all other harmonic resonances, solve
a scalar periodic rate equation, and finally identify its signed
linear response with a positive prescribed evolution.

\subsection{Physical parameter dependence and harmonic selection}

All parameter perturbations in this section use the original physical
variables on the unchanged half-strip. In particular the floor does
not move when $m$ or $a$ changes.

\begin{lemma}[Analyticity at positive noise]\label{lem:parameter-analyticity}
Locally at every $a>0$ under consideration, the completed inverse
$R_{i\theta}B:\HH_\sigma\to\EE_\sigma$ is holomorphic in $(m,a)$,
uniformly in $\theta\in\R$ on a sufficiently small complex parameter
neighborhood. The stationary triple $(M,M_y,M_F)$ and $N_0=\ell M$
are locally analytic. The response \eqref{eq:physical-response} is
locally analytic in $(\lambda,c,m,a)$ near every imaginary frequency.
The selected crossing functions $\cH,\OmH$ are real analytic on
$U=(0,a_0)\times(2,3)$.
\end{lemma}

\begin{proof}
At a fixed base $(m,a)$, regard the density, derivative, and outgoing
trace as separate components of $\EE_\sigma$. A perturbed inverse solves
\begin{equation}\label{eq:parameter-inverse}
 h=R_{i\theta}B[g-\delta_m h+\delta_a h_y].
\end{equation}
The right side uses the first and second components of $h$. The
uniform triple bound makes its linear perturbation small on the
ambient product for small complex $\delta_m,\delta_a$, independently
of $\theta$. Its Neumann series is analytic. Every output has its
derivative and trace identified by the completed weak equation;
closedness of the distributional derivative and passage in that
equation retain these identities in the limit. Its floor flux is
the perturbed natural flux. For real parameters, homogeneous
uniqueness identifies it with the actual prescribed inverse.

Likewise the nearby normalized stationary density is $M+h$, where
\begin{equation}\label{eq:parameter-stationary}
 h=R_0B[-\delta_m(M+h)+\delta_a(M_y+h_y)],\qquad \int_Dh=0.
\end{equation}
The same series gives analyticity. For real parameters normalized
stationary uniqueness identifies this density with the stationary
state supplied by Proposition~\ref{prop:frozen}, including its
positivity. Weighted norm closeness alone is not used for that
conclusion. Local frequency analyticity follows from the ordinary
resolvent identity as in Proposition~\ref{prop:divergence}. This
proves the response statement.

The simple zero and the nonzero real coupling derivative in
\eqref{eq:crossing} and \eqref{eq:response-root} give analytic crossing
functions by the implicit function argument. Local solutions patch
on $U$ by uniqueness of the selected branch and its crossing.
No uniform radius in $m$ as $a\downarrow0$, and no convergence of
parameter derivatives in that limit, is asserted.
\end{proof}

\begin{lemma}[A base point with only first-harmonic resonance]\label{lem:selection}
There is $(a,m_*)\in U$ such that, writing
\[
 c_*=\cH(a,m_*),\qquad \Omega_*=\OmH(a,m_*),\qquad T=2\pi/\Omega_*,
\]
one has
\begin{equation}\label{eq:nonresonance}
 \partial_m\OmH(a,m_*)\ne0,\qquad
 1-\resp(ik\Omega_*;m_*,a,c_*)\ne0
                         \quad(k\in\mathbb Z\setminus\{-1,1\}).
\end{equation}
The first-harmonic zeros are simple.
\end{lemma}

\begin{proof}
For each fixed $k\ge2$, Proposition~\ref{prop:return-root} gives, at
$c=\cH$,
\[
 \operatorname{Re}\lambda_k
 =a\left[-k^2D+\frac D{\gamma_1}\gamma_k\right]+O(a^{3/2})<0
\]
for small $a$ depending on this $k$. The strict sign follows from
$\gamma_k/k^2<\gamma_1$. The real analytic complex-valued function
\[
 r_k(a,m)=1-\resp(ik\OmH(a,m);m,a,\cH(a,m))
\]
is not identically zero on $U$. Indeed, for this fixed $k$ take $a$
small enough that $ik\OmH$ lies in its fixed return-root disk.
A zero of $r_k$ would, by the imaginary-axis matching in
\eqref{eq:response-match}, give an eigenvalue-one equation for $K$
there. Its only root is $\lambda_k$, which has negative real part,
a contradiction.

Also $\OmH(a,m)\to\omega(m)$ uniformly on $I$, while
\[
 \tau_m=-\int_0^L[m-b(x)]^{-2}\dd x<0.
\]
Choose two interior means with different $\omega$. For all sufficiently
small $a$, $\OmH(a,\cdot)$ has different values at those means.
Thus $\partial_m\OmH$ is not the zero analytic function on $U$.
The zero set of a nonzero real analytic function on this connected
open set has empty interior: vanishing on a ball makes all its
power-series coefficients zero and propagates along a chain of
overlapping balls. For complex-valued $r_k$, apply this observation
to a nonzero real or imaginary part.

The zero sets of all $r_k$ and of $\partial_m\OmH$ are closed with
empty interior in $U$. Successively choose nested closed balls,
each inside the preceding interior, avoiding the next zero set,
with radii tending to zero and the first ball compactly inside $U$.
Completeness gives a point avoiding every set. The static multiplier
is nonzero by \eqref{eq:static-response}; negative harmonics are
complex conjugates. This proves \eqref{eq:nonresonance}.
Simplicity at the first harmonics follows from
\eqref{eq:response-root}.
\end{proof}

The selection is existential and uses only a fixed-$k$ asymptotic
for each integer $k$. It does not require a small-noise expansion
uniform over infinitely many harmonics.

\subsection{The periodic inverse and the rate equation}

From now on fix the point and the period in Lemma~\ref{lem:selection}.
Use Fourier coefficients
$u_k=T^{-1}\int_0^T u(t)e^{-ik\Omega_*t}\dd t$, so that
\[
 \norm u_{H^1(0,T;E)}^2
   =T\sum_{k\in\mathbb Z}(1+k^2\Omega_*^2)\norm{u_k}_E^2
\]
for a Hilbert space $E$.

\begin{lemma}[High-frequency response and periodic inverse]\label{lem:periodic-inverse}
At fixed positive parameters,
\begin{equation}\label{eq:response-tail}
 \resp(i\theta;m,a,c)\longrightarrow0\quad(|\theta|\to\infty).
\end{equation}
The coefficientwise map
\[
 (\mathcal R_TBg)_k=\mathcal R_{k\Omega_*}g_k
\]
is bounded from $H^1_{\per}(0,T;\HH_\sigma)$ to
$H^1_{\per}(0,T;\EE_\sigma)$, with its density, gradient, trace, and
floor identities. The linear rate-response multiplier on
$H^1_{\per}(0,T)$ is compact.
\end{lemma}

\begin{proof}
For fixed $h\in\HH_{\sigma,0}$, $t\mapsto S(t)h$ is Bochner
integrable by the frozen gap. Approximate it in $L^1_t\HH_\sigma$
by finitely many step functions and integrate their scalar
exponentials. Formula \eqref{eq:ordinary-resolvent} then shows
$R_{i\theta}h\to0$ in $\HH_\sigma$. Apply
\eqref{eq:domain-energy} to this resolvent:
\[
 \norm{(R_{i\theta}h)_F}_{\OO_\sigma}^2
 \le C_a\norm{R_{i\theta}h}_{\HH_\sigma}^2+
              C_a\norm h_{\HH_\sigma}\norm{R_{i\theta}h}_{\HH_\sigma}.
\]
Its outgoing trace tends to zero as well. Approximate $g_c$ in
$\HH_\sigma$ by compact smooth $g_j$ supported away from the floor.
For fixed $j$, $Bg_j$ is an ordinary zero-mass datum, so the preceding
argument applies. The uniform completed trace estimate controls
the approximation error independently of $\theta$, proving
\eqref{eq:response-tail}.

Parseval's identity and \eqref{eq:triple-bound} give the stated
periodic bound first on finite Fourier sums, then by completion.
The completed weak equation identifies the limiting derivative
and traces, preserves mass zero, and retains the natural floor input.
The density belongs to $C_{\per}(0,T;\HH_\sigma)$.
Finally the scalar multipliers tend to zero by
\eqref{eq:response-tail}, so finite Fourier truncations approximate
the linear response in operator norm on the same $H^1$ space.
\end{proof}

Keep $a$ fixed and allow $c,m$ near $c_*,m_*$. For a small real
$n\in H^1_{\per}(0,T)$ solve the completed equation
\begin{equation}\label{eq:nonlinear-response}
 h_n=R_TB\{n g_c+n[-c h_n+c^2(h_n)_y]\},\qquad \int_Dh_n(t)=0.
\end{equation}
Here and below $R_TB$ denotes the density component when the entire
triple is understood. The one-dimensional Sobolev product inequality
and Lemma~\ref{lem:periodic-inverse} bound the last term on
$H^1_t\EE_\sigma$ by $C\norm n_{H^1}$ times the triple norm.
A Neumann series therefore gives a unique small analytic solution.
Parameter analyticity follows uniformly in Fourier frequency from
\eqref{eq:parameter-inverse}. Put $F_n=M+h_n$. Its weak equation
has precisely the prescribed mean $m+cn$ and diffusion $a+c^2n$,
with the voltage traces and the natural floor flux. At this stage
$F_n$ is allowed to be signed.

The actual self-consistency equation is the scalar equation
\begin{equation}\label{eq:rate-equation}
 \mathcal F(n,c,m):=n-\ell h_n=0 .
\end{equation}
Its linearization at $n=0$ is the Fourier multiplier
\[
 d_k=1-\resp(ik\Omega_*;m,a,c).
\]
At the base point only $k=\pm1$ are singular. Its restriction to
the complement is boundedly invertible: finitely many nonzero
$d_k$ are bounded away from zero, and
\eqref{eq:response-tail} gives $|d_k|\ge1/2$ in the remaining tail.

\begin{proposition}[A nonconstant self-consistent rate]\label{prop:rate-branch}
For all sufficiently small $r>0$ there are real analytic parameter
values $c(r^2),m(r^2)$ near $c_*,m_*$ and real
$n_r\in H^1_{\per}(0,T)$ solving \eqref{eq:rate-equation}, with exact
first Fourier coefficient $(n_r)_1=r$ and
$\norm{n_r}_{H^1}\to0$.
\end{proposition}

\begin{proof}
Write
$n=ze^{i\Omega_*t}+\overline z e^{-i\Omega_*t}+w$, with
$w_{\pm1}=0$. The complementary equation has an analytic solution
$w=w(z,\overline z,c,m)$ by its bounded inverse and the local
implicit function argument. Equivalently, subtract its linear part
and use the resulting contraction. Time translations commute with
the equation and its projections, so uniqueness makes $w$
translation covariant. In the remaining first coefficient,
a real analytic monomial $z^p\overline z^{\,q}$ must have $p-q=1$.
Consequently
\begin{equation}\label{eq:reduced-equation}
 \mathcal F_1=z\,g(|z|^2,c,m),\qquad
 g(0,c,m)=1-\resp(i\Omega_*;m,a,c).
\end{equation}
The latter equality holds also for nearby parameters: the linear
rate operator remains diagonal in Fourier modes, so its first-mode
input has zero complementary projection. The persistent
complementary inverse thus gives $w_z(0)=0$.

Let $\lambda(c,m)$ be the simple characteristic root near
$i\Omega_*$. Along its crossing curve,
\[
 \lambda_m+\lambda_c\,\partial_m\cH=i\,\partial_m\OmH,\qquad
 \det_{\R}(\lambda_c,\lambda_m)
       =\operatorname{Re}\lambda_c\,\partial_m\OmH\ne0.
\]
The determinant uses the real and imaginary parts as its two rows.
If $d(\lambda,c,m)=1-\resp(\lambda;m,a,c)$, differentiation of
$d(\lambda(c,m),c,m)=0$ multiplies these two parameter columns
by the same nonzero complex number $-d_\lambda$. It therefore
preserves nonsingularity. The real Jacobian of $g(0,c,m)$ in
$(c,m)$ is invertible. The finite-dimensional analytic implicit
function argument now gives $c(\rho),m(\rho)$ with
$g(\rho,c(\rho),m(\rho))=0$ for small real $\rho$.
Take $\rho=r^2$, $z=r>0$. This proves the proposition.
The use of two free physical parameters avoids any assumption
about a cubic coefficient in a one-parameter bifurcation.
\end{proof}

\subsection{Positivity and the physical solution}

An $H^1_t\HH_\sigma$ perturbation of $M$ need not be pointwise
nonnegative. The next argument uses a positive period map instead.

\begin{lemma}[Comparison of prescribed evolutions]\label{lem:flow-comparison}
On each fixed finite interval $[0,t_*]$, let $S_n(t,0)$ be the
prescribed positive conservative evolution with mean $m+cn(t)$
and diffusion $a+c^2n(t)$. For small real $n$ and nearby $m,c$,
\begin{equation}\label{eq:flow-comparison}
 \norm{S_n(t_*,0)-S_*(t_*)}_{\HH_\sigma\to\HH_\sigma}
 \le C_{a,t_*}\bigl(|m-m_*|+\norm n_\infty\bigr),
\end{equation}
where $S_*$ is the frozen evolution at $(m_*,a)$.
\end{lemma}

\begin{proof}
For the same initial datum $f$ put $u_n=S_n(\cdot,0)f$.
The homogeneous forward energy is uniform on a small positive
coefficient neighborhood, so
$\norm{u_n}_{L^\infty_t\HH_\sigma}
 +\norm{(u_n)_y}_{L^2_t\HH_\sigma}\le C_{a,t_*}\norm f_{\HH_\sigma}$.
The difference $u_n-S_*f$ solves the base forced equation with
zero initial datum and input
\[
 g=-[(m-m_*)+cn]u_n+c^2n(u_n)_y .
\]
Its $L^2_t\HH_\sigma$ norm is at most the right side of
\eqref{eq:flow-comparison} times $\norm f$. The completed weak
flux identity and prescribed uniqueness identify this difference
with the solution in Lemma~\ref{lem:causal}. Estimate
\eqref{eq:causal-bound} proves the assertion. This uses a causal
forcing estimate, not an assumed forced renormalization formula.
\end{proof}

\begin{proof}[Proof of Theorem~\ref{thm:main}]
Choose the base point in Lemma~\ref{lem:selection}. By the frozen gap,
there is an integer $p$ with
$\norm{S_*(pT)}_{\HH_{\sigma,0}\to\HH_{\sigma,0}}<1/4$.
Keep this integer fixed. For the solution in
Proposition~\ref{prop:rate-branch}, choose $r>0$ so small that
\eqref{eq:flow-comparison} at $pT$ is less than $1/4$.
The prescribed $pT$ map then contracts every mass-zero difference
by a factor less than $1/2$.

Iterate that map from the normalized nonnegative density $M$.
Every iterate is nonnegative and has mass one. The affine
mass-one space in $\HH_\sigma$ is complete, so the iterates
converge to its unique fixed point $F_0^*$. Positivity and mass
one pass to the limit, since their defining sets are closed.
If $P_n=S_n(T,0)$, periodic coefficients and uniqueness give
$S_n(pT,0)=P_n^p$. The operator $P_n$ commutes with this power,
and $P_nF_0^*$ is another mass-one fixed point of it.
Uniqueness gives $P_nF_0^*=F_0^*$.
Thus $F^*(t)=S_n(t,0)F_0^*$ is nonnegative and $T$-periodic.

The signed response $F_n=M+h_n$ in \eqref{eq:nonlinear-response}
lies in $C_t\HH_\sigma$ with its derivative and traces in the
prescribed energy class. Its completed weak equation is the
same prescribed homogeneous equation. Prescribed uniqueness
identifies it with its evolution from $F_n(0)$. Since it is
$T$-periodic and of mass one, $F_n(0)$ is a fixed point of the
$pT$ map on the full affine space, and hence equals $F_0^*$.
This proves $F_n=F^*\ge0$ without a lower bound for $M$.

Set, with $m=m(r^2)$ and $c=c(r^2)$,
\begin{equation}\label{eq:physical-reconstruction}
 y_*=m-cN_0(m,a),\qquad a_*=a-c^2N_0(m,a).
\end{equation}
At the base point $c=O(a)$ and $N_0$ is uniformly bounded, while
$m>1$. First choose $a$ small and then $r$ small. This makes
$y_*,a_*,c$ positive. The exact rate equation gives
$N=N_0+n_r$, so its first Fourier coefficient is $r>0$.
Moreover
\[
 y_*+cN=m+cn_r,\qquad a_*+c^2N=a+c^2n_r .
\]
The prescribed equation for $F_n$ is therefore the original
nonlinear equation with both feedback terms. The choices
$v_E=2,v_F=1,y_L=1$ satisfy all remaining physical parameter
inequalities.

We verify the requested solution properties. The density belongs
to $H^1_{\per}(0,T;\HH_\sigma)$, hence to
$C_{\per}(0,T;\HH_\sigma)\subset C_{\per}(0,T;L^1(D))$.
The continuous representative has mass one at every time.
Exponential-weight Cauchy--Schwarz bounds all finite conductance
moments uniformly in time, in particular $\int y^2F$.
The positive entropy part is bounded by $\int(F^2+F)$.
For the negative part choose the normalized density
$q_0(v,y)=C e^{-y^2}$. The pointwise inequality
$s\log(s/q_0)\ge s-q_0$, together with the second moment and the
finite positive part, bounds $\int|F\log F|$ uniformly over
the period, with $0\log0=0$.

The low-conductance zero traces and high-conductance flux reset
are the strong completed voltage traces in
\eqref{eq:nonlinear-response}, equivalently those of the
prescribed evolution. Its floor trace is zero for
\[
 Q=(m+cn_r-y)F-(a+c^2n_r)F_y
      =(y_*+cN-y)F-(a_*+c^2N)F_y.
\]
Finally pair $Q$ with any smooth periodic time test and any
smooth test on the closed voltage interval. This pairing lies
in $L^1(\dd y)$ by the weighted energy bounds. The PDE,
integration by parts in $t,v$, and the integrable outgoing and
reset traces put its distributional $y$ derivative in $L^1$
on every tail. It has an absolutely continuous tail
representative with a finite limit at infinity. Its own
$L^1$ integrability forces that limit to be zero. This is the
vanishing distributional flux trace in the theorem.
All required boundary, entropy, continuity, and moment
properties follow for the nonnegative full density $F=F_n$.
\end{proof}


\appendix
\section{Process identification and the reflecting boundary}\label{app:boundary}

This appendix supplies the boundary facts used in Sections
\ref{sec:linear}--\ref{sec:returns}. The interior regularity theorem
is used only after a distributional equation has been established.
The reflecting condition is constructed separately.

\subsection{From the forward density to the reset process}

\begin{proof}[Proof of Proposition~\ref{prop:process}]
Let $F_t=S(t)h$. Put
$\Phi(v)=(\cos(2\pi v),\sin(2\pi v))$ and lift $F_t$ to a
probability $\nu_t$ on $\R^3$: choose $(v,y)$ with density $F_t$
and an independent fair sign, then set $(\xi,z)=(\Phi(v),\pm y)$.
This probability is supported on the unit circle times $\R$ and
need not have an ambient density. Away from the circle cut
$p=(1,0)$ use the inverse angle $v(\xi)\in(0,1)$ and define
\[
 B_\xi=2\pi J(v(\xi),|z|)(-\xi_2,\xi_1),\qquad
 B_z=m\,\operatorname{sign}z-z,\qquad
 A=\operatorname{diag}(0,0,2a).
\]
Choose any finite Borel angular drift at the cut and extend it
by zero off the circle. The value at the cut does not affect
the forward equation, since $\nu_t$ gives it zero mass.

For a compact smooth ambient test $\Psi$, its even pullback is
\[
 \varphi(v,y)=\tfrac12[\Psi(\Phi(v),y)+\Psi(\Phi(v),-y)].
\]
It has equal voltage endpoint values and $\varphi_y(v,0)=0$.
The reset terms cancel in the forward Green identity and the
low-conductance traces vanish. Integration by parts in $y$ has
no remaining floor contribution. The result is precisely the
Euclidean Fokker--Planck equation with drift $B$ and covariance
$A$, whose second-order term is $(1/2)A:D^2=a\partial_{zz}$.
The lifted probabilities are narrowly continuous and
\[
 \int_0^t\int_{\R^3}(|B|+|A|)\dd\nu_s\dd s<\infty,
\]
since $|B|\le C(1+|z|)$ on their support and $F$ has a bounded
first conductance moment.

Trevisan's superposition theorem \cite[Definition 2.2,
condition (4), and Theorem 2.5]{Trevisan15} supplies a continuous
martingale solution with these one-time laws. The angular
component lies on the circle at all rational times, and hence
at all times by continuity. Localized coordinate and product
tests show that its angular martingale has zero quadratic
variation, while its signed conductance satisfies
\begin{equation}\label{eq:signed-conductance}
 \dd Z_t=(m\,\operatorname{sign}Z_t-Z_t)\dd t+\sqrt{2a}\dd W_t .
\end{equation}
Here $W$ is scalar Brownian motion in the full solution
filtration, by the martingale quadratic variation and the scalar
L\'evy characterization. The stochastic integral, maximal
inequality, localized It\^o formula, and scalar characterization
used here can, for example, be taken from
\cite[Theorems 2.32, 4.10, 4.17, 4.19--4.20 and the scalar case
of Theorem 5.11]{DeVecchi22}.

We justify folding directly. Stop at
$\tau_R=\inf\{t:|Z_t|\ge R\}$ and apply the localized It\^o formula
to $\varphi_\delta(z)=\sqrt{z^2+\delta^2}$, made bounded outside
$[-R-1,R+1]$. For $0<\delta<1$,
\[
 a\E\int_0^{t\wedge\tau_R}\varphi_\delta''(Z_s)\dd s
       \le2R+2+(m+R)t .
\]
On $|z|\le\delta$ its second derivative is at least
$(2\sqrt2\,\delta)^{-1}$. Hence
\[
 \E\int_0^{t\wedge\tau_R}\one_{\{|Z_s|\le\delta\}}\dd s
       \le C_{a,R,t}\delta .
\]
Let $\delta\downarrow0$ and then $R\uparrow\infty$ to conclude
that $Z$ spends zero time at zero. The drift terms in the
same It\^o formula converge to those for $|Z|$; the stochastic
integrals converge uniformly in probability by isometry and
the maximal inequality. Consequently
\[
 K_t=|Z_t|-|Z_0|-\int_0^t(m-|Z_s|)\dd s
       -\sqrt{2a}\int_0^t\operatorname{sign}(Z_s)\dd W_s
\]
is a continuous nondecreasing limit of
$a\int_0^t\varphi_\delta''(Z_s)\dd s$.
On a path interval where $|Z|\ge d>0$, the latter increments
are bounded by $at\delta^2/d^3$ and vanish. Thus $K$ increases
only when $|Z|=0$. The martingale
$W'_t=\int_0^t\operatorname{sign}(Z_s)\dd W_s$ has quadratic
variation $t$, so it is scalar Brownian motion in the same
filtration. The folded $Y=|Z|$ satisfies
\eqref{eq:reset-process}.

The half-line Skorokhod map
$x(t)\mapsto x(t)+\sup_{s\le t}(-x(s))_+$
has supremum-norm Lipschitz constant at most two. Its integral
equation with drift $m-Y$ is a contraction on any interval of
length less than $1/2$. Iteration proves pathwise uniqueness
and makes $Y$ a measurable function of $Y_0,W'$. Brownian
increments are independent of the initial full sigma-field.
If $p(y)=\int_0^1h(v,y)\dd v$, Fubini therefore shows that the
conditional density of $V_0$ given the entire future $Y$ path
is $h(v,Y_0)/p(Y_0)$ on the set $p(Y_0)>0$.

For a fixed continuous conductance path, before a reset,
\begin{equation}\label{eq:voltage-flow}
 V_t=A_tV_0+C_t,\quad
 A_t=\exp\!\left[-\int_0^t(Y_s+1)\dd s\right],\quad
 C_t=\int_0^t\exp\!\left[-\int_s^t(Y_r+1)\dd r\right]2Y_s\dd s .
\end{equation}
The $C^1$ function $H(t)=(1-C_t)/A_t$ satisfies
$H'(t)=-J_F(Y_t)/A_t$. A grazing threshold hit would make
$V_0$ a critical value of $H$. Critical values of a $C^1$
real function on a bounded interval have Lebesgue measure
zero: cover the zero-derivative set by open intervals where
$|H'|<\epsilon$ and bound the lengths of their images by
$\epsilon$ times their total length, then let $\epsilon$
decrease to zero. Exhaust the time intervals.
The conditional density just proved excludes such a hit.
A first hit from below has $J_F\ge0$, and thus in fact $Y>1$.

After a reset the initial voltage is zero, so a separate
point-start argument is required. Until a floor encounter the
auxiliary unreflected conductance and voltage flow have a joint
density at every positive time. Indeed It\^o's formula gives
their distributional forward equation. The smooth fields
$\partial_y$ and
$\partial_t+J\partial_v+(m-y)\partial_y$, and their bracket
$(2-v)\partial_v-\partial_y$, span all time--space directions
for $v<2$. H\"ormander's theorem
\cite[Theorem 1.1]{Hormander67} gives a smooth space--time
density there. The explicit flow stays below $2$; narrow
continuity against compact tests upgrades the almost-every-time
density identification to every positive time.

Starting after a reset with $Y>1$, wait $\delta=1/n$ before the
next hit or first floor encounter. The law at $\delta$ restricted
to this survival event is dominated by the auxiliary law and
has a joint density. Future Brownian increments are independent
of the past event. The conditional-density critical-value
argument therefore applies after $\delta$.
The countable collection of these events covers all possible
finite next hits. Independence after stopping times follows by
approximating them from above by dyadic stopping times and
using Brownian continuity and deterministic-time independence.
Thus the argument repeats after every reset.

Before its first cut encounter the circle path inverts to the
ordinary voltage equation. It cannot reach zero from inside,
since $V_t\ge V_0A_t>0$. At a non-grazing outgoing hit the
angular speeds on both sides of the cut are positive, as
$J_F>0$ and $J_0=2Y>0$. Continuity of $Y$ makes them bounded
away from zero in a neighborhood. The path crosses in the
positive direction and enters the unique voltage trajectory
from zero. It cannot stick on an interval; in addition its
total occupation of the cut is zero by the one-time marginals.
The pathwise identity for $X$ gives \eqref{eq:count-bound}.
Since $Y$ is continuous, there is no finite-time accumulation
of hits. The same bound and the first conductance moment make
the expected count finite. This identifies the circle process
with the reset process, proving its law assertion.

For smooth $q$ supported in $(0,t)\times(1,\infty)$ apply
piecewise It\^o to $\varphi(s,v,y)=vq(s,y)$, whose reset jump
is $q$. Comparing with the forward Green formula gives
\eqref{eq:count-measure} against these tests. The trace of a
nonnegative forward density is nonnegative in its kinetic
trace measure; this also follows by the bounded negative-part
renormalization in the same Green identity. Both measures are
therefore locally finite positive measures. All hits have
$Y>1$. A spike at a deterministic time would put the circle
path at the cut, which its marginal excludes. There is no
initial spike for a density. Exhaust the conductance and
time endpoints. Finite expected count controls the total
measure, and monotone convergence extends equality first to
bounded Borel tests and then to integrable tests. This proves
the complete statement.
\end{proof}

\subsection{Local backward equations, including the floor}

We give the local construction for killed values with the weights
of Section~\ref{sec:returns}. Their interior sources belong to the
smooth class specified before Lemma~\ref{lem:cycle-energy} and obey
the stated Gaussian weighted bound. We also treat bounded
full-process resolvents with compact smooth source.
In both cases write
\begin{equation}\label{eq:backward-source-equation}
 \G u=\lambda u-r .
\end{equation}
The source is zero for a killed harmonic. Nonzero killed sources
used below are the constant source in the static block and the
smooth corrected-mode residual. The full-process case has
$\operatorname{Re}\lambda>0$ and compact smooth bounded $r$.

First establish the interior distributional equation. On an interior
rectangle use the unreflected, unreset process until exit from a larger
rectangle. Its conductance is explicit and its voltage is
\eqref{eq:voltage-flow}; the stochastic flow is triangular in the
initial variables, with positive Jacobian $e^{-t}A_t$. For a compact
smooth initial density $\varphi$, Taylor expansion of its inverse flow
and inverse Jacobian gives
\begin{equation}\label{eq:local-transpose-expansion}
 P_h^T\varphi
  =\varphi+h[-\partial_v(J\varphi)
             -\partial_y((m-y)\varphi)+a\varphi_{yy}]+o(h)
\end{equation}
in local $L^1$ and in the fixed Gaussian-weighted $L^1$ used to pair
with $u$. To see the remainder explicitly, the conductance increment
has mean $(m-y)h+O(h^2)$, variance $2ah+O(h^2)$, and $j$th absolute
moment $O(h^{j/2})$; its random contribution to the voltage increment
is $O(h^{3/2})$ in every fixed finite moment. The inverse Jacobian
is $1+h(y+2)+O_{L^p}(h^{3/2})$. Bounded third derivatives of
$\varphi$ control the Taylor remainder. Outside an enlarged compact
set the transported density has a Gaussian small-time tail, which
dominates the fixed exponential-quadratic weight. The probability
of leaving the larger rectangle before $h$ is $O(e^{-c/h})$;
higher stopped weights make its contribution $o(h)$.
For a smooth killed source the reward increment is $hr+o(h)$ on
compact sets of starts. The same localization and the higher-weight
reward bounds justify this term in the paired limit.
Dynamic programming, pairing with $\varphi$ and division by $h$,
now give \eqref{eq:backward-source-equation} in distributions.

We use the real-coefficient regularity theorem literally at complex
frequency. Introduce a real variable $z$ and set
$U(z,v,y)=e^{-\lambda z}u(v,y)$. Then
$(\partial_z+\G)U=-e^{-\lambda z}r$.
The fields $\partial_y$ and
$\partial_z+J\partial_v+(m-y)\partial_y$, with the same bracket
$(2-v)\partial_v-\partial_y$, span all three directions.
Apply \cite[Theorem 1.1]{Hormander67} to the real and imaginary
parts of $U$ separately. Thus $u$ is interior smooth.

At $v=0,y>0$ extend the auxiliary voltage a fixed small distance
below zero. An interior physical path never exits through zero;
the high-conductance block argument still bounds the hitting
time from the added finite voltage interval. The weak equation
and interior regularity on this extension make $u(0,y)$ smooth
for $y>0$. For a constant reset jump, the outgoing datum is
therefore smooth on $y>1$. Its positive outgoing speed gives
the terminal trace. At $v=1,y<1$, choose a small rectangle
straddling that face on which $J\le-j_0<0$. Stop the auxiliary
process on its exit. Its exit time is bounded by the voltage
width divided by $j_0$. Retain the original data on reachable
physical sides and extend side data continuously to the added
part. A physical start cannot exit through the added upper
voltage side. The local value agrees with $u$ inside, solves
the weak equation across the face, and is smooth there.
The deterministic exit bound suffices for any fixed complex
discount; no global hitting estimate outside the physical
voltage interval is assumed.

Point starts at positive conductance have no grazing hit by the
positive-time density argument above. For a start at $y=0$,
before any finite threshold hit the conductance first reaches
$1/2$ strictly earlier; start the same argument at that stop.
Continuous dependence of the Skorokhod and voltage equations
then implies almost-sure continuity of $\tf,Y_\tf$ in the
initial state at a transverse hit. Before the hit the voltage
on each shorter compact interval stays a positive distance
below threshold, and the final transverse crossing persists.
The stopped $4\beta$ weights and stronger discount moment give
uniform integrability on compact sets of starts. Thus killed
values are continuous, including at the floor. For a bounded
full-process source, the same coupling on finite intervals
and the discounted tail prove continuity.

Fix $0<v_-<v_+<1$ and choose $y_c>0$ so that $q=-J$ satisfies
$0<q_-\le q\le q_+$ on the resulting floor rectangle.
The backward equation becomes the forward voltage-clock problem
\begin{equation}\label{eq:floor-clock}
 q u_v=au_{yy}+(m-y)u_y-\lambda u+r,\qquad
 u_y(v,0)=0 .
\end{equation}
Give it the continuous initial data at $v_-$ and the upper-side
data at $y_c$. After a smooth lifting of that side, use a mixed
Neumann--Dirichlet basis on $(0,y_c)$. Its $q$-weighted mass
matrix is uniformly positive. Testing the finite systems by
their solutions yields
\[
 \tfrac12\partial_v\int_0^{y_c}q|w|^2\dd y
        +a\int_0^{y_c}|w_y|^2\dd y
 \le C\int_0^{y_c}(|w|^2+|\text{source}|^2)\dd y .
\]
The drift floor term is $-m|w(v,0)|^2/2$, with favorable sign.
Weak Hilbert limits, the variational equation, and the
zero-data energy inequality give existence and uniqueness,
with continuous $L^2$ time representative. Testing with $w_v$
after a positive voltage interval bounds $w_v,w_{yy}$ in $L^2$;
repeated differentiated estimates give the local boundary
regularity needed for It\^o. The absence of a diffusion floor
term in the variational identity encodes its Neumann trace.
For continuous initial data, approximation and the maximum
bound give the same solution. The full Galerkin passage and
positive-time arguments are detailed for the analogous
Dirichlet clock in Appendix~\ref{app:clock}; here the variational
space imposes only the upper-side zero value.

Until this rectangle is exited, physical voltage decreases
strictly. In its inverse clock the conductance generator is
$q^{-1}[a\partial_{yy}+(m-y)\partial_y]$, its discount is
$\lambda/q$, and its source is $r/q$. It\^o applied to the
constructed solution cancels the drift and the Neumann
derivative cancels the regulator. Physical exit time is
bounded by $(v_+-v_-)/q_-$. Stopped dynamic programming has
the same boundary data, so this solution equals $u$.
This proves the actual floor derivative for $v>0$, without
applying an interior theorem at a boundary.

\subsection{Growth for a constant jump and a bounded-source energy}

We first supply the growth assertion used in
Proposition~\ref{prop:constant-jump}. Suppose the reset jump is $A$.
Choose $p>\max\{1,r_0\}$ and a fixed $y_c>m,1$ so large that
\[
 (\G_y+r_0)y^p\le-d(y+r_0L),\qquad
 \G_y=(m-y)\partial_y+a\partial_{yy},\quad y\ge y_c .
\]
This follows from
$\G_y y^p=-py^p+pmy^{p-1}+ap(p-1)y^{p-2}$.
For $S=\inf\{t:Y_t\le y_c\}$, stopped It\^o gives
\[
 \E e^{r_0S}+
 \E\int_0^S e^{r_0t}(Y_t+r_0L)\dd t\le C(1+y^p).
\]
Recall $X=\log[2/(2-V)]\in[0,L]$, $L=\log2$.
Between hits $\dot X=Y-(e^X-1)$, and each reset decreases $X$ by $L$.
Including the reset jumps in $e^{r_0t}X_t$ gives
\[
 L\sum_{t_j\le S}e^{r_0t_j}
       \le X_0+\int_0^S e^{r_0t}(Y_t+r_0L)\dd t .
\]
Patching the killed martingales until $S$ consequently yields
\[
 u(z)=\E_z[e^{-\lambda S}u(Z_S)]
                +A\E_z\sum_{t_j\le S}e^{-\lambda t_j}.
\]
To remove upper-conductance and time localizations before asserting
this equality, choose $\beta<\beta'<1/2$ and enlarge $y_c$ so that
$(\G_y+r_0+d_0)w_{\beta'}\le0$ above it. The upper-exit term
is bounded by a fixed starting-state constant times
$w_\beta((H-m)/\varepsilon)/w_{\beta'}((H-m)/\varepsilon)$,
which tends to zero. The time endpoint is bounded by
$e^{-d_0t}w_{\beta'}(\eta_0)$.
The displayed count bound controls the spike sum independently.
Since $u$ is bounded on $y=y_c$, this proves
\begin{equation}\label{eq:polynomial-growth}
 |u(v,y)|\le C_{a,m,\lambda}(\norm f_\beta+|A|)(1+y^p).
\end{equation}
A bounded interior source adds a controlled
$\E\int_0^S e^{r_0t}|r|\dd t$ and gives the same conclusion.
These are fixed-$a$ growth constants.

The independent dual estimate of Section~\ref{sec:linear} needs
only bounded full-process source values. We state it separately
to keep its proof independent of killed-value gradients.

\begin{proposition}[Bounded-source backward energy]\label{prop:backward-source}
For bounded compact smooth $r$ and $\lambda=s+i\theta$ with $s>0$,
let
\[
 u(z)=\E_z\int_0^\infty e^{-\lambda t}r(Z_t)\dd t .
\]
It is bounded by $\norm r_\infty/s$, has zero reset jump,
satisfies \eqref{eq:backward-source-equation}, and has the local
Neumann floor constructed above. For $s$ larger than a fixed
constant,
\begin{equation}\label{eq:backward-energy}
 (s-C)\norm u_{\HH_{-\sigma}}^2+
          a\norm{u_y}_{\HH_{-\sigma}}^2
 \le \norm r_{\HH_{-\sigma}}\norm u_{\HH_{-\sigma}},
\end{equation}
uniformly in $\theta$. For $s$ also larger than the forward
growth bound, this value is the transpose of the full forward
resolvent.
\end{proposition}

\begin{proof}
The local construction above applies to this bounded value,
including its homogeneous reset and actual floor derivative.
Let $\rho=e^{-2\sigma y}$. On an exhausted domain the real
energy identity has volume coefficient
\[
 C_\sigma(y)=(\tfrac12-\sigma)y+1+\sigma m+2a\sigma^2,
\]
gradient term $-a\int|u_y|^2\rho$, floor term
$-(m/2+a\sigma)\int|u(v,0)|^2\dd v$, and voltage term
$\frac12\int\rho[J_F|u_F|^2-J_0|u_0|^2]\dd y$.
The floor and low-conductance voltage terms are nonpositive.
At high conductance $u_F=u_0$ and $J_F<J_0$, so that term is
nonpositive too. Since $\sigma>1$, $C_\sigma$ is bounded above.

Here is a localization that proves the first global gradient
bound rather than assuming it. Remove rectangles of voltage
width $\delta^3$ and conductance width $\delta$ around $(1,1)$
and $(0,0)$. Choose cutoffs $\chi_\delta$ with
$|\partial_v\chi_\delta|\le C\delta^{-3}$,
$|\partial_y\chi_\delta|\le C\delta^{-1}$ and zero floor
derivative. Add a large-$y$ cutoff $\zeta_R$ and test with
$\chi_\delta^2\zeta_R^2\overline u\rho$.
Near a corner $u$ is bounded, $|J|=O(\delta)$, and area is
$O(\delta^4)$. The voltage-cutoff error is $O(\delta^2)$,
the squared conductance-cutoff error is $O(\delta^2)$, and the
drift-cutoff error is $O(\delta^3)$. Young's inequality absorbs
the unknown gradient cross terms before any global gradient
integrability is assumed. A reset-cutoff mismatch is confined
to an outgoing interval with $\int J_F\dd y=O(\delta^2)$.
At large $y$, boundedness of $u$ and the exponential weight
remove all tail errors after the same absorption.

The resulting gradient bounds are uniform in $\delta,R$,
so Fatou gives $u_y\in\HH_{-\sigma}$. Repeat the passage with
this known gradient: the first-derivative cutoff errors now
tend to zero by Cauchy--Schwarz, since
$\norm{(\chi_\delta)_yu}_2=O(\delta)$.
The exact energy inequality and the preceding signs yield
\eqref{eq:backward-energy}. Finally Proposition~\ref{prop:process}
and Fubini identify the value paired against any positive
initial density with the full forward Laplace resolvent.
Splitting signed data and density in the weighted spaces
give the claimed transpose identification.
\end{proof}

\section{The bounded clock comparison}\label{app:clock}

We prove \eqref{eq:clock-sup} with no regularity modulus assumed
for the terminal datum. This is the local estimate that permits
the operator norm, rather than strong-operator, limit in
Proposition~\ref{prop:return-norm}.

Return to the scaled conductance $\eta=(y-m)/\sqrt a$ and the OU
expression $\mathcal A=\partial_{\eta\eta}-\eta\partial_\eta$.
The deterministic phase interval is $[0,\tau]$, with
$\tau=s(\log2)$ and scalar coefficient
$Z(s)=[m-b(x(s))]^{-1}$ from \eqref{eq:phase}.
These phase objects determine the reversed coefficient
$Z(\tau-t)$ used below.

Let $I_R=(-R,R)$, $\dd\gamma=e^{-\eta^2/2}\dd\eta$,
$E=L^2(I_R,\dd\gamma)$, and $V=H^1_0(I_R)$.
For a smooth positive coefficient $q$ with fixed upper and
lower bounds consider
\begin{equation}\label{eq:clock-general}
 q u_t=(\mathcal A-\lambda)u,\qquad
 u(t,\pm R)=0,\qquad u(0)=f .
\end{equation}
The coefficient and its needed derivatives are bounded on
the fixed rectangle; $\lambda$ ranges in a compact set.

\subsection{Construction and positive-time estimates}

Choose a smooth Dirichlet basis with finite spans dense in
$V$ and $E$, and project the initial datum in $E$. On each
finite span the $q$-weighted mass matrix is uniformly positive,
so the finite linear system has a solution. Testing with $u_N$
and using $\langle\mathcal A u,v\rangle_\gamma
=-\langle u_\eta,v_\eta\rangle_\gamma$ gives
\[
 \frac{\dd}{\dd t}\int q|u_N|^2\dd\gamma
 =-2\norm{(u_N)_\eta}_2^2
   -2\operatorname{Re}\lambda\norm{u_N}_2^2
   +\int q_t|u_N|^2\dd\gamma .
\]
Consequently
\begin{equation}\label{eq:clock-energy}
 \sup_{0\le t\le\tau}\norm{u_N(t)}_2^2
            +\int_0^\tau\norm{(u_N)_\eta}_2^2\dd t
       \le C\norm f_2^2 .
\end{equation}
Take a weak subsequence in the Hilbert space $L^2_tV$.
The $L^\infty_tE$ bound survives: testing the lower norm
bound on every measurable time set $A$ gives
$\int_A\norm u_E^2\dd t\le C\norm f_2^2|A|$.
Integrate the finite equations by parts in time against
smooth tests with values in any fixed finite span. Every
term is linear, so weak convergence passes it, including
$q_t$ and the initial datum. Density gives the full
variational equation. Strong time compactness is unnecessary.

Multiplication by $q^{\pm1}$ preserves $V$, and the equation
gives $u_t\in L^2_tV'$. Time Steklov averages justify the
energy identity: on a smaller interval the averages lie in
$H^1_tV$, so differentiating their squared $E$ norm is
legitimate. The supremum norm of the difference of two
averages is bounded by its time-average $E$ norm squared
plus twice its $L^2_tV$ norm times its derivative's
$L^2_tV'$ norm. Translation continuity makes both terms
tend to zero. Forward averages near the initial endpoint
and backward averages near the final endpoint give a
$C_tE$ representative and the limiting energy identity.
Equivalently one may apply the same argument to the
$q$-weighted norm, retaining its bounded $q_t$ term.
The initial weak datum and the energy limsup give $u(0)=f$.
The zero-data energy inequality proves uniqueness. The
entire Galerkin sequence therefore has this limit.

Testing the finite equation by $(u_N)_t$ instead gives
\[
 \frac{\dd}{\dd t}\norm{(u_N)_\eta}_2^2
       +c\norm{(u_N)_t}_2^2\le C\norm{u_N}_2^2 .
\]
Choose an earlier time with bounded gradient using
\eqref{eq:clock-energy}, then integrate. Passing weakly
gives, for each $t_*>0$,
\begin{equation}\label{eq:clock-smoothing}
 \sup_{t_*\le t\le\tau}\norm{u(t)}_V^2
       +\int_{t_*}^\tau\norm{u_t}_2^2\dd t
       \le C_{t_*}\norm f_2^2 .
\end{equation}
The equation gives $u_{\eta\eta}\in L^2$ there.
Repeated differentiated equations, with the smooth $q$
commutators, and the one-dimensional equation give the
local $C^1_tC^2_\eta$ representative on positive-time
closed subrectangles. The $V$ bound extends to each
positive time by weak $V$ compactness and strong $E$
continuity; hence the spatial output is continuous.
No compatibility at the initial Dirichlet corners is used.

For $q=1$ the equation is autonomous. Its derivative $v=u_t$
satisfies the same homogeneous equation. Choose a time
strictly before $\tau/2$ with controlled $\norm v_2$ using
\eqref{eq:clock-smoothing}, then apply the energy and
positive-time estimates to $v$. The one-dimensional
Dirichlet inequality now gives
\begin{equation}\label{eq:clock-autonomous-derivative}
 \sup_{\tau/2\le t\le\tau}
       \norm{(\mathcal A-\lambda)u_0(t)}_\infty
       \le C\norm f_2 .
\end{equation}

The maximum comparison is asserted for the weak limit,
not for pointwise values of a Galerkin approximation.
Regularize the modulus by $\sqrt{|u|^2+\delta^2}$ and test
the positive part above a scalar exponential barrier.
The extra second-derivative term has nonnegative sign
and can be dropped. The $V$ energy permits the limit
$\delta\downarrow0$. This proves the maximum bound
and its version with an inhomogeneous bounded source.

\subsection{Perturbation in supremum norm}

Take $q_\varepsilon=1+\varepsilon\eta\chi(t)$, where
$\chi(t)=Z(\tau-t)$ and
$\varepsilon R\norm\chi_\infty\le1/2$.
Put $w=u_\varepsilon-u_0$ and $h=1-q_\varepsilon$.
Its exact equation is
\[
 q_\varepsilon w_t=(\mathcal A-\lambda)w+
                          h(\mathcal A-\lambda)u_0,\qquad w(0)=0 .
\]
The forcing integration by parts reads
\[
 \langle h\mathcal A u_0,w\rangle_\gamma
   =-\int (h\overline{w_\eta}+h_\eta\overline w)
                       (u_0)_\eta\dd\gamma .
\]
Both $h,h_\eta$ are $O(\varepsilon)$ on the fixed rectangle.
Young's inequality and \eqref{eq:clock-energy} give
\begin{equation}\label{eq:clock-difference}
 \sup_{t\le\tau}\norm{w(t)}_2^2+
           \int_0^\tau\norm{w_\eta}_2^2\dd t
       \le C\varepsilon^2\norm f_2^2 .
\end{equation}
Choose $t_1\in[\tau/2,3\tau/4]$ with
$\norm{w_\eta(t_1)}_2\le C\varepsilon\norm f_2$.
The Dirichlet inequality on the bounded interval implies
$\norm{w(t_1)}_\infty\le C_R\varepsilon\norm f_2$.
On the remaining interval the divided equation has source
\[
 F=q_\varepsilon^{-1}h(\mathcal A-\lambda)u_0,\qquad
 \norm F_\infty\le C\varepsilon\norm f_2
\]
by \eqref{eq:clock-autonomous-derivative}. Modulus comparison
with $M'=C_\Lambda M+\norm F_\infty$ therefore gives
\[
 \norm{w(\tau)}_\infty
 \le e^{C_\Lambda\tau}\norm{w(t_1)}_\infty+
       \int_{t_1}^\tau e^{C_\Lambda(\tau-t)}\norm{F(t)}_\infty\dd t
 \le C\varepsilon\norm f_2
 \le C_R\varepsilon\norm f_\infty .
\]
This is \eqref{eq:clock-sup}. Rescaling the compact family
of intervals $\tau(m)$ to a unit interval makes its
constants uniform for $m\in I$.

\subsection{Identification with the stopped return}

On the bounded conductance interval the physical phase speed
is positive. Reverse the remaining phase time as in
\eqref{eq:clock-equation}. The inverse-clock conductance has
generator $q_\varepsilon^{-1}\mathcal A$ and killing rate
$\lambda/q_\varepsilon$. Apply stopped It\^o to the
constructed positive-time classical representative.
Its equation cancels the drift. The clock bounds make
physical duration at most $\tau/q_-$, and the maximum
bound controls the limiting stopped values.

There is no simultaneous terminal-phase and side exit.
To verify this, extend $q_\varepsilon$ smoothly to a
globally positive bounded coefficient. The unrestricted
inverse-clock OU law satisfies its smooth uniformly
parabolic forward equation. Interior regularity followed
by narrow continuity gives a density at each positive
phase time. A side exit exactly at the fixed terminal
phase would put this conductance at $R$ or $-R$,
an event of probability zero. Thus continuous $f$ is
attained on surviving terminal paths and the zero
Dirichlet value on earlier side exits. Dominated
convergence in the stopped identity identifies the
weak evolution with the killed return kernel.
Approximation in $E$ and the maximum bound give the
same identification for arbitrary continuous initial
data used in the comparison.

For the floor clock \eqref{eq:floor-clock}, replace $V$
by $\{w\in H^1(0,y_c):w(y_c)=0\}$ and use Lebesgue
measure. The $q$-weighted mass remains positive, and
the first energy has the favorable floor drift
$-m|w(0)|^2/2$. In the derivative estimate bound the
drift term against $w_v$ by Young's inequality before
any spatial integration by parts. The same weak
passage, uniqueness, and positive-time argument apply.
Reflection contributes $u_y\dd K=0$ in It\^o.
The corresponding nondegenerate auxiliary clock law
also excludes a tie of the absorbing upper side
and initial voltage. This completes the construction
used for the reflecting boundary in Appendix~\ref{app:boundary}.


\section*{AI use statement}
OpenAI Codex was used extensively for mathematical exploration,
proof development, internal consistency checks, literature retrieval,
and preparation and revision of this manuscript. These internal AI
checks do not constitute independent proof verification.
Two additional Codex contexts provided reader and referee feedback
on the initial draft. This version incorporates one combined
revision responding to that feedback.
The author is responsible for the final manuscript.

\begin{thebibliography}{99}

\bibitem{CDZ22}
J.~A. Carrillo, X.~Dou, and Z.~Zhou,
\emph{A simplified voltage-conductance kinetic model for interacting
neurons and its asymptotic limit},
preprint (2022),
\href{https://arxiv.org/abs/2203.02746v1}{arXiv:2203.02746v1}.

\bibitem{CR25}
J.~A. Carrillo and P.~Roux,
\emph{Nonlinear partial differential equations in neuroscience:
from modelling to mathematical theory},
preprint (2025),
\href{https://arxiv.org/abs/2501.06015v1}{arXiv:2501.06015v1}.

\bibitem{DeVecchi22}
F.~De Vecchi,
\emph{Foundations in Stochastic Analysis},
lecture notes, WS 2021/22, January 31, 2022.
\href{https://mgubi.github.io/store/iam/stochastic-analysis-ss22/Foundations_in_Stochastic_Analysis.pdf}{Author's lecture notes}.

\bibitem{FM24}
C.~Fonte Sanchez and S.~Mischler,
\emph{On the Voltage-Conductance kinetic equation},
preprint (2024),
\href{https://arxiv.org/abs/2408.02471v1}{arXiv:2408.02471v1}.

\bibitem{FMS26}
C.~Fonte Sanchez, S.~Mischler, and D.~Salort,
\emph{Existence of solutions to the Voltage-Conductance kinetic
equation in a general conductivity regime},
preprint (2026),
\href{https://arxiv.org/abs/2607.22112v1}{arXiv:2607.22112v1}.

\bibitem{Hormander67}
L.~H\"ormander,
\emph{Hypoelliptic second order differential equations},
Acta Math. \textbf{119} (1967), 147--171.
\href{https://archive.ymsc.tsinghua.edu.cn/pacm_download/117/6044-11511_2006_Article_BF02392081.pdf}{Primary scan}.

\bibitem{Tao11}
T.~Tao,
\emph{An Introduction to Measure Theory},
Graduate Studies in Mathematics, vol.~126, American Mathematical
Society, 2011. Author's preliminary version, 2012.
\href{https://terrytao.wordpress.com/wp-content/uploads/2012/12/gsm-126-tao5-measure-book.pdf}{Author's preliminary text}.

\bibitem{Trevisan15}
D.~Trevisan,
\emph{Well-posedness of Multidimensional Diffusion Processes with
Weakly Differentiable Coefficients},
preprint (2015),
\href{https://arxiv.org/abs/1507.01357v2}{arXiv:1507.01357v2}.

\end{thebibliography}

\end{document}
